\documentclass[11pt,letterpaper]{amsart}
\usepackage[T1]{fontenc}
\usepackage{lmodern}
\usepackage[margin=1.1in]{geometry}
\usepackage{amsmath,amssymb,amsthm,mathtools}
\usepackage{microtype,booktabs,enumitem,needspace}
\usepackage{xcolor}
\usepackage{hyperref}
\hypersetup{colorlinks=true,linkcolor=blue!45!black,citecolor=blue!45!black,urlcolor=blue!45!black,pdftitle={Hyperrigidity in low dimensions and a dimension-33 counterexample},pdfauthor={Caleb Barnett}}
\setlist[enumerate]{label=\textup{(\roman*)},leftmargin=2em}
\newtheorem{theorem}{Theorem}[section]
\newtheorem{lemma}[theorem]{Lemma}
\newtheorem{proposition}[theorem]{Proposition}
\newtheorem{corollary}[theorem]{Corollary}
\newtheorem{maintheorem}{Theorem}
\renewcommand{\themaintheorem}{\Alph{maintheorem}}
\theoremstyle{definition}
\newtheorem{definition}[theorem]{Definition}
\theoremstyle{remark}
\newtheorem{remark}[theorem]{Remark}
\newtheorem{example}[theorem]{Example}
\numberwithin{equation}{section}
\DeclareMathOperator{\ex}{ex}
\DeclareMathOperator{\dist}{dist}
\DeclareMathOperator{\aff}{aff}
\DeclareMathOperator{\supp}{supp}
\DeclareMathOperator{\relint}{relint}
\DeclareMathOperator{\diam}{diam}
\DeclareMathOperator{\rank}{rank}
\DeclareMathOperator{\ran}{ran}
\DeclareMathOperator{\diag}{diag}
\newcommand{\R}{\mathbb R}
\newcommand{\C}{\mathbb C}
\newcommand{\N}{\mathbb N}
\newcommand{\ind}{\mathbf 1}
\newcommand{\B}{\mathcal B}
\title[Hyperrigidity in low dimensions]{Hyperrigidity in low dimensions\\and a dimension-33 counterexample}
\author{Caleb Barnett}
\date{Version 2, 23 October 2026}
\dedicatory{\normalfont\small\textbf{Working draft --- Version 2}\\
Not published or released as a preprint.}
\thanks{Prepared with substantial AI assistance and automated adversarial review. Independent human verification and a complete novelty assessment remain pending; no formal verification is claimed.}
\keywords{Hyperrigidity, unique extension property, compact convex sets, positive operator-valued measures, Riesz energy.}
\begin{document}
\begin{abstract}
We study the least affine dimension of a commutative counterexample to
Arveson's hyperrigidity conjecture. Using an inverse-square energy
estimate for supporting affine functions, we prove hyperrigidity for
every compact convex subset of $\R^3$ with closed extreme boundary.
The key step replaces part of a dilation while preserving supporting
normals, after which Borel interpolation on exposed segments avoids
disintegration. We also modify the Walsh dilation of Bilich--Scherer to
construct a counterexample in affine dimension $33$. Thus the least
affine dimension lies between $4$ and $33$. Further results exclude
several smaller construction mechanisms, including rank-one total
quadratic defects, Hardy compressions of multiplicity at most three,
and small polynomial perturbations of quadratic graphs. A finite-matrix
criterion and connections with matrix convexity and quantum channels
identify remaining directions toward dimension four.
\end{abstract}
\maketitle

\section{Introduction}
\label{sec:introduction}

Let $K\subset\R^d$ be a compact convex set with closed extreme boundary
$X=\ex K$. We write $A(K)$ for the complex vector space of continuous
affine functions on $K$ and regard
\[
                         A(K)|_X\subset C(X)
\]
as a function system. A unital representation $\pi:C(X)\to B(H)$ has
the \emph{unique extension property} (UEP) if the only unital completely
positive (UCP) map $\Phi:C(X)\to B(H)$ agreeing with $\pi$ on $A(K)|_X$ is
$\pi$ itself. By Arveson's characterization, the function system is
hyperrigid exactly when every unital representation has this property
\cite{Arv}. We use this formulation throughout.

Each point evaluation at an extreme point has the unique extension
property: a probability measure with that barycenter is a point mass.
The question is whether the same uniqueness survives for diffuse
representations. Bilich and Scherer showed that it need not
\cite{BS}. Their construction raises a geometric question: how small
can the affine dimension of a commutative counterexample be?

The present paper proves universal hyperrigidity in dimension three, gives an explicit
counterexample in dimension $33$, and structural restrictions on
attempts to bridge the gap. Its first main result is the following.

\Needspace{8\baselineskip}
\begin{maintheorem}\label{intro:positive}
Let $K\subset\R^3$ be a compact convex set with closed extreme
boundary. Then $A(K)|_{\ex K}$ is hyperrigid in $C(\ex K)$.
\end{maintheorem}

No smoothness or exposedness hypothesis is imposed. The theorem extends
the positive results in the first version of this paper. The new proof
retains every exceptional source point: countably many simple pieces
are handled separately, and the residual points are endpoints of exposed
segments. A dilation replacement preserves their supporting normals,
allowing bounded Borel interpolation between the segment endpoints.

The common mechanism is an energy estimate. If $n_x$ and $n_y$ are unit
outward normals at $x,y\in\partial K$, then
\begin{equation}\label{intro:eq:lift}
 0\le\langle n_x,x-y\rangle
 \le\frac14\bigl|(x+n_x)-(y+n_y)\bigr|^2.
\end{equation}
The normal lift $x\mapsto x+n_x$ takes values on
$\partial(K+\overline B)$, where $\overline B$ is the closed unit ball.
In dimension three, every nonzero measure on this parallel surface has
infinite inverse-square energy. On the other hand, a nonzero separated
block of an agreeing positive operator-valued measure would produce a
finite-energy measure. This contradiction proves the required
separation. The replacement step assigns each new target copy its
source normal, even when several sources share a nonsmooth target.
The complete proof is given in Theorem~\ref{thm:three}.

Our second main result is constructive.

\begin{maintheorem}\label{intro:example}
There is a compact convex body $K\subset\R^{33}$ with closed extreme
boundary $X$, a unital representation $\pi:C(X)\to B(H)$, and a unital
completely positive map $\Phi:C(X)\to B(H)$ such that
\[
                   \Phi|_{A(K)|_X}=\pi|_{A(K)|_X},
                   \qquad \Phi\ne\pi.
\]
The Hilbert space $H$ is separable and infinite dimensional, and the
spectral measure of $\pi$ has no point atoms. Moreover
$\dim_{\C}A(K)=34$.
\end{maintheorem}

The construction is a change of distribution in the Walsh dilation of
Bilich--Scherer, rather than a new dilation mechanism. Theorem~\ref{thm:main}
gives the complete data. The operator identities are proved before the
numerical parameters are chosen; the remaining containment estimate
is a scalar inequality for every infinite digit sequence. A continuous
barycentric kernel then gives the extension, and a quadratic function
provides an explicit positive witness of disagreement.

For comparison, \cite[Corollary 3.3]{BS} proves a construction in
ambient dimension $129$, states an optimization to dimension $57$, and
mentions an unverified dimension-$11$ proposal in Remark~3.4.
Scherer's planar theorem \cite[Theorem 3.8]{planar} excludes dimensions
at most two. Together, Theorems~\ref{intro:positive} and
\ref{intro:example} give
\begin{equation}\label{eq:dimension-range}
                            4\le d_{\min}\le33,
\end{equation}
where $d_{\min}$ denotes the least affine dimension under the
closed-extreme-boundary hypothesis. Whether the lower bound four is attained remains open. The positive theorem for compact spectrahedra
\cite[Theorem 3.7]{spectra} is a further point of comparison; it applies
to finite self-adjoint affine matrix pencils, not arbitrary projections
of spectrahedra.

The construction obstructions explain why simply changing its visible
geometry does not close this gap. A finite-dimensional signing lemma
rules out the quadratic-distance certificate for every stationary
three-point distribution and every geometric-norm vector encoding in
horizontal dimension at most four, under the nonnegative-potential and
exact-height assumptions of Theorem~\ref{enc:universal}. Thus even arbitrary directions cannot
make that certificate work in affine dimension four or five.
A stronger geometric argument rules out \emph{every finite strictly
convex graph} for the same stationary mechanism in affine dimension
four (Corollary~\ref{prox:cor:walsh}). It uses positive distance levels
of the represented set and requires no differentiability of the graph.
These are obstructions to a specified construction family, not universal
lower bounds on $d_{\min}$.

The general supporting-gap criterion, Theorem~\ref{bar:thm}, is useful
beyond these applications. Appendix~\ref{app:capacity} gives a
source in a four-dimensional convex body whose normal lift supports a
finite-energy measure, but whose representations still have the unique extension
property by a sharper metric. Appendices~\ref{app:heterogeneous}
and~\ref{app:correlated} treat varying and correlated distributions.
Appendix~\ref{app:direct} records two restrictions on direct operator
constructions. The later sections add independent construction obstructions and a
finite-POVM criterion, followed by a literature-guided discussion of
dimension four. Exploratory numerical searches and superseded
certificates are not needed for any proof in this paper.

\section{Representations, measures, and defects}
\label{sec:preliminaries}

All Hilbert spaces are complex. Representations and completely positive
maps are unital unless stated otherwise. A convex body has nonempty
interior in its ambient Euclidean space. We write $\overline B$ for the
closed unit ball and use $|\cdot|$ for Euclidean norms. For a compact
convex set $K$, its support function is
$h_K(u)=\max_{x\in K}\langle u,x\rangle$. An exposed point is the
unique minimizer of an affine function on $K$. A smooth boundary point has a unique unit outward
normal. A strictly convex body has no nondegenerate segment in its
boundary; every boundary point is then exposed.

Whenever $X=\ex K$ is compact, the coordinate functions and the constant
function generate $C(X)$. If $\aff K=\R^d$, their complex linear span
is $A(K)|_X$ and has dimension $d+1$. The dimension of a representation
space is a separate issue.

For a representation $\pi:C(X)\to B(H)$ we write $P$ for its spectral
measure. A UCP map $\Phi$ has a positive operator-valued measure $E$.
One may construct its bounded Borel extension from a Stinespring
representation $\Phi(f)=V^*\rho(f)V$: if $Q$ is the spectral measure of
$\rho$, set $E(B)=V^*Q(B)V$. All operator-valued measures below are
countably additive in the strong operator topology. Saying that a
representation is concentrated on a Borel set $Z$ means $P(Z)=I$;
this is stronger than an assertion about the closure of that set.

\begin{lemma}[Coordinate defects]\label{lem:defect}
Suppose $\Phi=V^*\rho(\,\cdot\,)V$ agrees with $\pi$ on the affine
coordinates $x_1,\ldots,x_d$. Then
\[
 D_j:=\Phi(x_j^2)-\pi(x_j)^2
   =\bigl[(I-VV^*)\rho(x_j)V\bigr]^*
      \bigl[(I-VV^*)\rho(x_j)V\bigr]\ge0.
\]
All $D_j$ vanish if and only if $VH$ reduces $\rho(C(X))$, and this
is equivalent to $\Phi=\pi$.
\end{lemma}
\begin{proof}
Expand the square and use $V^*\rho(x_j)V=\pi(x_j)$.
If it vanishes, $VH$ is invariant for the self-adjoint operator
$\rho(x_j)$ and hence reduces it. The coordinates generate $C(X)$,
so the compression is a representation agreeing with $\pi$ on a
generating set. The converse follows from the same identity.
\end{proof}

\begin{lemma}[Atomic representations]\label{lem:atomic}
Every Hilbert direct sum of point evaluations at points of $X$,
with arbitrary multiplicities, has the unique extension property.
\end{lemma}
\begin{proof}
Write $H=\bigoplus_x H_x$, with $\pi(f)|_{H_x}=f(x)I$.
For a unit vector $h\in H_x$, the scalar measure representing
$f\mapsto\langle\Phi(f)h,h\rangle$ has barycenter $x$ and is therefore
$\delta_x$. Indeed, a probability measure with an extreme barycenter
cannot be split into two parts with different conditional barycenters.
In a Stinespring representation this gives
\[
 \|\rho(f)Vh-f(x)Vh\|^2
       =\langle\Phi(|f-f(x)|^2)h,h\rangle=0.
\]
Thus $\rho(f)V=V\pi(f)$ on each summand and then on their Hilbert direct
sum. This proves equality, including all off-diagonal coefficients.
\end{proof}

In particular, finite-dimensional representations cannot witness a
counterexample. We shall also need a stronger local consequence of
extremality.

\begin{lemma}[Local vanishing of the defect]\label{het:lem:local}
Let $\Phi$ agree with $\pi$ on $A(K)|_X$, and let $D_j$ be the defects
above. If $x$ belongs to the spectral support of $\pi$, then
\[
 \lim_{r\downarrow0}
 \left\|P(B(x,r))\left(\sum_jD_j\right)P(B(x,r))\right\|=0.
\]
Consequently no coordinate defect, nor their sum, is bounded below by
a positive multiple of the identity.
\end{lemma}
\begin{proof}
Any sequence of probability measures on $X$ whose barycenters tend to
$x$ converges weakly to $\delta_x$: every weak limit has the extreme
barycenter $x$. By compactness, their second moments about $x$ tend
to zero uniformly as their barycenters approach $x$.
For a unit vector $h\in P(B(x,r))H$, the scalar measure of $\Phi$
at $h$ has barycenter within $r$ of $x$. Affine agreement gives
\[
 \left\langle\sum_jD_jh,h\right\rangle
 =\int_X|y-x|^2\,d\mu_h(y)
   -\left\langle\pi(|\,\cdot-x|^2)h,h\right\rangle
 \le\int_X|y-x|^2\,d\mu_h(y).
\]
The right side tends to zero uniformly in these unit vectors.
Positivity gives the norm estimate. Every such spectral neighborhood
is nonzero, which proves the last assertion.
\end{proof}

\section{A quadratic affine-barrier criterion}
\label{paper:pos:barriers}

The positive results in this section rest on a common observation.
An affine supporting function can bound the amount of a dilation
that passes from one part of the spectrum to another. If these
bounds vanish quadratically in a suitable metric, a nonzero
off-diagonal effect would produce a scalar measure of finite
inverse-square energy. A critical covering estimate rules out
such a measure. The metric need not be the ambient Euclidean
metric, and this flexibility will be useful below.

We retain the notation of Section~\ref{sec:preliminaries}:
$P$ is the spectral measure of $\pi$, and $E$ is the POVM of
an affine-agreeing UCP map $\Phi$. Their bounded Borel extensions
will be used without changing notation.

\begin{theorem}[Quadratic affine barriers]\label{bar:thm}
Let $\pi:C(X)\to B(H)$ be a unital representation, and let
$Z\subset X$ be Borel with $P(Z)=I$. Suppose that, for each
$x\in Z$, an affine function $b_x$ on $K$ is specified so that
its coefficients depend Borel measurably on $x$ and
\[
 b_x(x)=0,\qquad b_x(y)>0\quad(y\in K\setminus\{x\}).
\]
Suppose also that there are a compact metric space $(T,d_T)$,
a Borel map $\theta:Z\to T$, and a finite constant $M$ such that
\[
 b_x(y)\le M\,d_T(\theta(x),\theta(y))^2
 \qquad(x,y\in Z).
\]
Assume that every nonzero finite positive Borel measure $\mu$
on $Z$ satisfies
\[
 \iint_{Z\times Z}
 d_T(\theta(x),\theta(y))^{-2}\,d\mu(x)d\mu(y)=+\infty,
\]
where the integrand is $+\infty$ when the distance is zero.
Then $\pi$ has the unique extension property relative to
$A(K)|_X$.
\end{theorem}

We first establish the analytic facts needed for the proof.
For a finite positive measure $\nu$ on a metric space, write
\[
 I_2(\nu)=\iint d(u,v)^{-2}\,d\nu(u)d\nu(v).
\]
A Borel set has \emph{zero inverse-square capacity} here if
every nonzero finite positive measure concentrated on it has
infinite $I_2$-energy. The formulation in
Theorem~\ref{bar:thm} is the corresponding statement for the
pushforwards under $\theta$; it does not require a separate
measurability assertion about the image of $\theta$.

\subsection{Energy and localization}

\begin{lemma}[Rank-one domination]\label{lb:lem:rankone}
Let $P$ be a spectral measure on a compact metric space, let
$w$ be a bounded nonnegative Borel function, and suppose
\[
 |\eta\rangle\langle\eta|\le M\int w\,dP.
\]
Then, for $\mu_\eta(D)=\langle P(D)\eta,\eta\rangle$,
\[
                  \int w^{-1}\,d\mu_\eta\le M,
\]
with $w^{-1}=+\infty$ on the zero set of $w$.
\end{lemma}
\begin{proof}
Put $W=\int w\,dP$. For $\epsilon>0$, set
$h=(W+\epsilon I)^{-1}\eta$ and
$t=\langle\eta,h\rangle\ge0$. The assumed domination gives
\[
 t^2\le M\langle Wh,h\rangle
       \le M\langle(W+\epsilon I)h,h\rangle=Mt.
\]
Thus $t\le M$. The spectral theorem and monotone convergence
as $\epsilon\downarrow0$ prove the assertion.
\end{proof}

\begin{lemma}[Critical covering bound]\label{lb:lem:energy}
Let $T$ be a compact metric space. Suppose that, for every
sufficiently small $\epsilon>0$, it can be covered by at most
$M\epsilon^{-2}$ sets of diameter at most $\epsilon$.
Then $I_2(\nu)=+\infty$ for every nonzero finite positive
Borel measure $\nu$ on $T$.
In particular, this conclusion holds for the boundary of
every compact convex body in $\mathbb R^3$.
\end{lemma}
\begin{proof}
The covering sets may be replaced by their closures and then
refined to a Borel partition
$T=\bigsqcup_{j=1}^{N_\epsilon}T_j$, with
$N_\epsilon\le M\epsilon^{-2}$ and
$\operatorname{diam}T_j\le\epsilon$.
If $m=\nu(T)>0$, Cauchy--Schwarz gives
\[
 (\nu\times\nu)\{d(u,v)\le\epsilon\}
 \ge\sum_j\nu(T_j)^2
 \ge\frac{m^2}{N_\epsilon}
 \ge\frac{m^2}{M}\epsilon^2.
\]
The layer-cake formula and Tonelli's theorem therefore imply
\[
 I_2(\nu)
 \ge2\int_0^{\epsilon_0}t^{-3}
       (\nu\times\nu)\{d(u,v)\le t\}\,dt
 \ge\frac{2m^2}{M}\int_0^{\epsilon_0}\frac{dt}{t}
 =+\infty.
\]

For the final assertion, translate the convex body so that
$r\overline B_3\subset K\subset R\overline B_3$ for some
$r,R>0$. Its Minkowski functional $j_K$ is $1/r$-Lipschitz and
is at least $1/R$ on $S^2$. Hence the radial map
$u\mapsto u/j_K(u)$ is Lipschitz from $S^2$ onto $\partial K$.
The sphere admits the required $O(\epsilon^{-2})$ covers,
as follows from finitely many Lipschitz coordinate charts.
The radial map transfers this covering bound to $\partial K$.
\end{proof}

The next lemma isolates the operator argument. Unlike
Theorem~\ref{bar:thm}, it does not require each affine function
to expose a single point.

\begin{lemma}[Local annihilation]\label{paper:pos:localization}
Suppose $\Phi$ and $\pi$ agree on $A(K)|_X$. Let
$F,G\subset X$ be Borel. For each $x\in F$, let $b_x$ be
an affine function nonnegative on $K$. Suppose that
\[
 b_x(y)\ge\delta>0\quad(x\in F,\ y\in G),
 \qquad
 b_x(y)\le M\,d_T(\theta(x),\theta(y))^2
       \quad(x,y\in F),
\]
where $\theta:F\to T$ is Borel and $T$ is compact metric.
If every nonzero finite positive measure on $F$ has infinite
energy in the displayed metric, then $E(G)P(F)=0$.
\end{lemma}
\begin{proof}
Set $A=P(F)E(G)P(F)$. For each fixed $x\in F$, positivity
and affine agreement imply
\[
 0\le A
 \le\delta^{-1}P(F)\Phi(b_x)P(F)
 =\delta^{-1}P(F)\pi(b_x)P(F)
 \le\frac M\delta\int_F
           d_T(\theta(x),\theta(y))^2\,dP(y).
\]
If $A\ne0$, choose a unit vector $h$ such that
$\eta=A^{1/2}h\ne0$. Then
$|\eta\rangle\langle\eta|\le A$ and $\eta\in P(F)H$.
Lemma~\ref{lb:lem:rankone}, applied on $P(F)H$, gives
\[
 \int_Fd_T(\theta(x),\theta(y))^{-2}\,d\mu_\eta(y)
 \le M/\delta\qquad(x\in F).
\]
The same vector $\eta$ works for every $x\in F$. Tonelli's theorem
therefore gives
\[
 \iint_{F\times F}d_T(\theta(x),\theta(y))^{-2}
       \,d\mu_\eta(y)d\mu_\eta(x)
 \le (M/\delta)\|\eta\|^2<\infty,
\]
contrary to the energy hypothesis. Thus $A=0$.
Since $A=(E(G)^{1/2}P(F))^*(E(G)^{1/2}P(F))$, it follows that
$E(G)P(F)=0$.
\end{proof}

We will also use the following elementary passage from compact
sets to the full spectral measure.

\begin{lemma}\label{paper:pos:measure-identification}
Let $P$ be a normalized spectral measure and $E$ a normalized POVM
on a compact metric space. If $E(B)P(A)=0$ whenever $A,B$ are
disjoint compact sets, then $E=P$.
\end{lemma}
\begin{proof}
For an open set $O$, choose compact sets $A_j\uparrow O$.
Applying the hypothesis in both orders to $A_j$ and the
compact set $X\setminus O$, and taking strong limits, gives
\[
 E(X\setminus O)P(O)=0,\qquad
 E(O)P(X\setminus O)=0.
\]
The two effects sum to $I$, so $E(O)P(O)=P(O)$ and
$E(O)P(X\setminus O)=0$. Hence $E(O)=P(O)$.
The scalar matrix-coefficient measures are regular and agree
on open sets; they therefore agree on all Borel sets.
\end{proof}

\begin{proof}[Proof of Theorem~\ref{bar:thm}]
Let $\Phi$ be an agreeing UCP map. Fix disjoint nonempty
compact sets $A,B\subset X$, and put
\[
 F_m=\left\{x\in A\cap Z:
                    \min_{y\in B}b_x(y)\ge1/m\right\}.
\]
The minimum over $B$ is a continuous function of the affine
coefficients, so $F_m$ is Borel. Strict positivity away from
$x$ and compactness of $B$ show that $F_m\uparrow A\cap Z$.
Lemma~\ref{paper:pos:localization} gives $E(B)P(F_m)=0$.
Since $P(Z)=I$, strong convergence yields $E(B)P(A)=0$.
Lemma~\ref{paper:pos:measure-identification} now gives
$E=P$, and hence $\Phi=\pi$.
\end{proof}

\subsection{A graph consequence}

The criterion applies to represented subsets, independently of
the size of the rest of the extreme boundary. The following
consequence is useful when a construction begins with a
convex graph.

\begin{corollary}\label{cap:cor:graph}
Let $B\subset\mathbb R^D$ be compact and convex, and let $g$
be strictly convex and $C^{1,1}$ on an open convex neighborhood
of $B$. Put
\[
 K=\{(a,z):a\in B,\ g(a)\le z\le h\},
 \qquad h>\max_B g,
\]
and assume $\ex K$ is closed. Let $S\subset B$ be compact
and admit covers by $O(\epsilon^{-2})$ sets of diameter at
most $\epsilon$. Every representation concentrated on
\[
                 Z=\{(s,g(s)):s\in S\}
\]
has the unique extension property.
\end{corollary}
\begin{proof}
For $x=(s,g(s))$, take the affine barrier
\[
 b_x(a,z)=z-g(s)-\langle\nabla g(s),a-s\rangle.
\]
Strict convexity makes this nonnegative on $K$, with equality
only at $x$. A Lipschitz bound $L$ for $\nabla g$ on $B$
gives
\[
 b_{(s,g(s))}(t,g(t))\le\tfrac L2|s-t|^2
 \qquad(s,t\in S).
\]
Use $\theta(s,g(s))=s$. Lemma~\ref{lb:lem:energy} and
Theorem~\ref{bar:thm} give the conclusion.
\end{proof}

\section{Normal lifts and local separation}
\label{paper:pos:normal-lifts}

Supporting normals provide a canonical metric for the barrier
criterion. The useful feature is that they can be added to the
boundary points: this places the resulting metric space on the
boundary of a parallel convex body. In dimension three, that
surface has precisely the critical covering bound of
Lemma~\ref{lb:lem:energy}.

For $x\in K$, let
\[
 N_K(x)=\{u\in\mathbb R^d:
               \langle u,y-x\rangle\le0\text{ for every }y\in K\}
\]
be the outward normal cone. For a unit vector $n_x\in N_K(x)$,
put $z_x=x+n_x$. The point $z_x$ lies on
$\partial(K+\overline B_d)$, and its metric projection onto
$K$ is $x$. Indeed, for $y\in K$,
\[
 |x+n_x-y|^2
 =1+|x-y|^2+2\langle n_x,x-y\rangle\ge1,
\]
with equality only when $y=x$. For any two such choices,
\begin{equation}\label{lb:eq:lift}
 0\le\langle n_x,x-y\rangle
 \le\langle n_x-n_y,x-y\rangle
 \le\tfrac14|z_x-z_y|^2.
\end{equation}
The middle inequality uses the supporting property at $y$;
the last follows from $4\langle a,b\rangle\le|a+b|^2$.

\begin{corollary}[Normal-lift criterion]\label{cap:thm}
Let $\pi:C(X)\to B(H)$ have spectral measure concentrated
on a Borel set $Z\subset X$. Suppose a Borel unit normal
$n(x)$ exposes $x$ for every $x\in Z$. If every nonzero finite
positive Borel measure concentrated on
\[
                        L=\{x+n(x):x\in Z\}
\]
has infinite inverse-square energy, then $\pi$ has the unique
extension property.
\end{corollary}
\begin{proof}
The map $x\mapsto x+n(x)$ is Borel and injective. Its image
is Borel as well: if $p_K$ is the continuous metric projection
onto $K$, then
\[
 L=\{z\in\partial(K+\overline B_d):
        p_K(z)\in Z,\ z-p_K(z)=n(p_K(z))\}.
\]
Apply Theorem~\ref{bar:thm} with
$b_x(y)=\langle n(x),x-y\rangle$ and
$\theta(x)=x+n(x)$. The quadratic bound is
\eqref{lb:eq:lift}, and the energy assumption passes to
pushforward measures.
\end{proof}

We need a local version in which the normal need not expose
the source point. This allows flat supporting faces in the general
three-dimensional argument of the next section.

\begin{lemma}[Separated normal selections]\label{lb:lem:sep}
Let $K\subset\mathbb R^3$ be a compact convex body, and suppose
$\Phi$ and $\pi$ agree on $A(K)|_X$. Let $F,G\subset X$ be
Borel. If there is a Borel unit normal selection
$x\mapsto n_x$ on $F$ such that
\[
             \langle n_x,x-y\rangle\ge\delta>0
             \qquad(x\in F,\ y\in G),
\]
then $Q(G)VP(F)=0$ and $E(G)P(F)=0$.
\end{lemma}
\begin{proof}
Use the affine functions
$b_x(y)=\langle n_x,x-y\rangle$ in
Lemma~\ref{paper:pos:localization}. By \eqref{lb:eq:lift},
their values on $F$ are bounded by one quarter of the squared
distance between the normal lifts. Every nonzero pushforward
measure on these lifts has infinite energy, since the lifts
lie on $\partial(K+\overline B_3)$ and
Lemma~\ref{lb:lem:energy} applies to that surface. The conclusion
$E(G)P(F)=0$ also gives $Q(G)VP(F)=0$ because the square of its
norm on a vector is the corresponding quadratic form of $P(F)E(G)P(F)$.
\end{proof}

For exposed points, a separating selection can be made on
successively larger compact sets. We record the elementary
selection argument to make the measurability explicit.

\begin{lemma}[Compact-graph selection]\label{lb:lem:selection}
Let $F,T$ be compact metric spaces, and let
$\Gamma\subset F\times T$ be compact with every fiber
$\Gamma_x$ nonempty. There is a Borel map $s:F\to T$ with
$(x,s(x))\in\Gamma$ for every $x\in F$.
\end{lemma}
\begin{proof}
Fix a countable dense sequence in $T$. At step $j$, choose
the first closed ball of radius $2^{-j}$ centered in that
sequence which meets the current nonempty compact fiber, and
intersect the fiber with that ball. These choices are Borel.
Indeed, for each fixed finite list of earlier choices, the
nonemptiness test is the projection of the original compact
graph intersected with finitely many fixed closed balls;
this projection is compact. The possible earlier lists form
a countable Borel partition, so induction proves the claim.
The remaining fibers are nested nonempty compact sets with
diameters tending to zero. Their unique intersection point
is the limit of the selected centers, and hence depends
Borel measurably on $x$.
\end{proof}

\section{Hyperrigidity in affine dimension three}
\label{sec:dimension-three}

\begin{theorem}\label{thm:three}
Let $K\subset\mathbb R^3$ be a compact convex body and assume that
$X=\ex K$ is closed. Every unital representation of $C(X)$ has the
unique extension property relative to $A(K)|_X$. Consequently,
$A(K)|_X$ is hyperrigid in $C(X)$.
\end{theorem}

Write $F_u=\{y\in K:\langle u,y\rangle=h_K(u)\}$ for $u\in S^2$.
We retain the normal-cone notation $N_K(x)$ from the preceding section.
A boundary point is smooth when $N_K(x)\cap S^2$ is a singleton.
Throughout, $\pi:C(X)\to B(H)$ is a representation and
$\Phi(f)=V^*\rho(f)V$ is a UCP map agreeing with $\pi$ on affine
functions. Here $V:H\to L$ is an isometry and $\rho:C(X)\to B(L)$ is a
unital representation. Their spectral measures are $P,Q$, and
$E(B)=V^*Q(B)V$. Bounded Borel functional calculus is understood.

\begin{lemma}[Separable reduction]\label{lem:separable}
To prove uniqueness for an affine-agreeing pair $(\pi,\Phi)$, or an
identity $Q(B)VP(A)=0$, it suffices to treat separable source and
Stinespring spaces. Borel spectral concentration assumptions persist.
\end{lemma}
\begin{proof}
Fix a vector witnessing failure of the asserted identity. The C*-algebra
generated by $\pi(C(X))$ and $\Phi(C(X))$ is separable because $C(X)$
is separable. The closed orbit of this vector under that algebra is a
separable reducing subspace $H_0$ for both maps. The space
$L_0=\overline{\operatorname{span}\rho(C(X))VH_0}$ is separable, reduces $\rho$, and
contains $VH_0$; it supplies the restricted Stinespring representation.
Reducing subspaces for a representation also reduce its bounded Borel
functional calculus. Hence $P(A)H_0\subset H_0$ and $Q(B)L_0\subset L_0$,
so either type of witness is retained. The same observation preserves
all Borel spectral concentration assumptions.
\end{proof}

\subsection{Contact localization}

\begin{lemma}[Borel contact localization]\label{lem:contact}
Let $S\subset X$ be the smooth locus and $n:S\to S^2$ its unique
normal map. If Borel sets $A\subset S$ and $B\subset X$ satisfy
\[
                  \langle n(x),x-y\rangle>0
                    \qquad(x\in A,\ y\in B),
\]
then $Q(B)VP(A)=0$.
\end{lemma}
\begin{proof}
By Lemma~\ref{lem:separable}, assume the representation spaces are
separable. A positive summable combination of scalar spectral measures
from countable dense sets of unit vectors gives a finite scalar measure $\nu$ with respect to which every scalar spectral measure
of $P,Q$ is absolutely continuous. Increasing compact subsets of $A$ and
$B$ exhaust each set modulo a $\nu$-null set. The normal map is continuous
relative to $S$: the unit-normal correspondence has closed graph and
compact values, so a limit of unique normals at a smooth limiting point
must be its unique normal. On each product of these compact sets the
displayed gap has a positive minimum. Lemma~\ref{lb:lem:sep} applies, and
successive strong limits give the claim.
\end{proof}

\subsection{The source decomposition}

\begin{lemma}\label{lem:countable}
For a compact convex body in $\mathbb R^3$, the nonexposed nonsmooth
extreme points form a countable set. Its two-dimensional exposed faces
also form a countable collection.
\end{lemma}
\begin{proof}
Let $x$ be nonexposed and nonsmooth. Choose linearly independent
$u_1,\ldots,u_k\in N_K(x)$ spanning $\operatorname{span}N_K(x)$ and put
$u=(\sum_i u_i)/|\sum_i u_i|$. The cone is pointed, so the
denominator is nonzero.
Nonnegativity of the supporting losses implies
\[
F_u=\bigcap_{v\in N_K(x)}\{y\in K:\langle v,x-y\rangle=0\}.
\]
This face has dimension at most $3-k$. Nonsmoothness gives $k\ge2$,
whereas nonexposedness gives $\dim F_u\ge1$. Hence $k=2$ and $F_u$
is an exposed segment with $x$ an endpoint. For an interior point $p$
of that segment, $N_K(p)=N_K(x)$: all normals at $x$ vanish on the segment,
and a normal at an interior point supports both endpoints.

There are at most countably many exposed segments $F$ having a
two-dimensional normal cone in their relative interiors. To see this,
choose a nondegenerate closed subsegment $F_0\subset\relint F$ and a
nondegenerate closed arc $C_0\subset N_K(F)\cap S^2$, where
$N_K(F)$ denotes the common normal cone on $\relint F$. The patch
$F_0+C_0$ lies on $\partial(K+\overline B_3)$ and has positive area:
the segment direction and the unit-speed circular-arc direction are
orthogonal, so its area is the product of their lengths. Different
segments give disjoint patches, since metric projection onto $K$ sends
each patch into the relative interior of its own segment. The parallel
surface has finite area, so only countably many disjoint positive-area
patches are possible. Each segment has two endpoints.

Finally, the relative interiors of distinct two-dimensional exposed faces
are disjoint positive-area Borel subsets of $\partial K$, which also has
finite area. Thus these faces are countable.
\end{proof}

\begin{lemma}\label{lem:measurable}
The smooth locus $S$ and the exposed-point set $\mathcal E$ in $X$ are
Borel. Let $\{F_j\}$ list the two-dimensional exposed faces, and put
\[
 Z=X\setminus\left(\mathcal E\cup C\cup\bigcup_j F_j\right),
 \qquad C=\{\text{nonexposed nonsmooth extreme points}\}.
\]
Then $Z$ is Borel, every $x\in Z$ is smooth, and
\[
                     F_{n(x)}=[x,\tau(x)]
\]
for a unique other endpoint $\tau(x)\in X\setminus\{x\}$. The map
$\tau:Z\to X$ is Borel.
\end{lemma}
\begin{proof}
The compact-valued unit-normal correspondence has closed graph, so its
diameter is upper semicontinuous. Its singleton locus $S$ is Borel.
The face correspondence $u\mapsto F_u$ has compact graph, and
$u\mapsto\diam F_u$ is upper semicontinuous. Thus the smooth exposed
points are exactly the Borel set
$\{x\in S:\diam F_{n(x)}=0\}$. The remaining exposed points are
$(X\setminus S)\setminus C$, a Borel set because $C$ is countable by
Lemma~\ref{lem:countable}.

For $x\in Z$, its supporting face cannot have dimension zero or two.
It is therefore a segment, with $x$ one endpoint because $x$ is extreme.
The other endpoint is extreme in that face and hence in $K$.
For completeness, measurability of this endpoint needs no selection
theorem. If $D\subset K$ is compact, define
\[
m_D(x)=\max\{|y-x|:y\in F_{n(x)}\cap D\},
\]
with value $-\infty$ for an empty intersection. Closedness of the face
graph and compactness of $D$ show that $m_D$ is upper semicontinuous
relative to $Z$. The point $\tau(x)$ is the unique maximizer defining
$m_K(x)$. For an open set $O\subset K$,
\[
                  \tau(x)\in O
          \quad\Longleftrightarrow\quad m_{K\setminus O}(x)<m_K(x).
\]
Both functions are Borel; hence $\tau$ is Borel.
\end{proof}

\begin{lemma}[Borel endpoints of exposed segments]\label{lem:segment-endpoints}
Let $K\subset\mathbb R^3$ be a compact convex body. The set
\[
 U=\{u\in S^2:F_u\text{ is a nondegenerate segment}\}
\]
is Borel. There are Borel maps $a,b:U\to\ex K$ such that
$F_u=[a(u),b(u)]$ and $a(u)\ne b(u)$ for every $u\in U$.
\end{lemma}
\begin{proof}
The face correspondence has compact graph. Consequently its diameter
and the maximal area of a triangle with vertices in $F_u$ are upper
semicontinuous functions of $u$. The set $U$ is characterized by
positive diameter and zero maximal triangle area, so it is Borel.

For $1\le j\le3$, put
\[
 \alpha_j(u)=\min_{x\in F_u}x_j,
 \qquad \beta_j(u)=\max_{x\in F_u}x_j.
\]
The functions $\alpha_j$ are lower semicontinuous and $\beta_j$ are
upper semicontinuous. Partition $U$ into the Borel sets $U_j$ on which
$j$ is the first coordinate satisfying $\beta_j(u)>\alpha_j(u)$.
For $u\in U_j$, let $a(u)$ and $b(u)$ be the unique points of $F_u$
minimizing and maximizing that coordinate. They are its two endpoints,
so they belong to $\ex K$.

Here is a direct measurability check, requiring no selection theorem.
For compact $D\subset K$, the extended-valued function
\[
 \alpha_{D,j}(u)=\min\{x_j:x\in F_u\cap D\},
 \qquad \min\varnothing=+\infty,
\]
is lower semicontinuous. Indeed, along any convergent sequence of
normals whose corresponding minima have a finite lower limit, compactness
provides a convergent subsequence of minimizers, and the face graph is
closed. For any relatively open $O\subset K$ and $u\in U_j$,
\[
 a(u)\in O
 \quad\Longleftrightarrow\quad
 \alpha_{K\setminus O,j}(u)>\alpha_j(u).
\]
The equivalence uses uniqueness of the minimizing endpoint and
compactness of $F_u\cap(K\setminus O)$. The right-hand side is Borel.
Thus $a$ is Borel on each $U_j$. Applying the same argument to maxima,
with $\max\varnothing=-\infty$, proves the assertion for $b$.
\end{proof}

\begin{lemma}\label{lem:good}
Every representation concentrated on $\mathcal E$, on a countable
subset of $X$, or on $X\cap F$ for a single proper exposed face $F$
has the unique extension property. A countable Borel partition into
sets with this property also has this property.
\end{lemma}
\begin{proof}
By Lemma~\ref{lem:separable}, it is enough to work in separable
spaces, where compact approximation uses a common finite scalar control
measure as above. For exposed sources, first fix disjoint compact sets $A\subset\mathcal E$
and $B\subset X$. For each positive integer $m$, consider the compact graph
\[
\Gamma_m=\{(x,u)\in A\times S^2:
 u\in N_K(x),\ \min_{y\in B}\langle u,x-y\rangle\ge1/m\}.
\]
Its compact projection $A_m$ increases to $A$. A Borel selection of its nonempty compact fibers is supplied by Lemma~\ref{lb:lem:selection}. Lemma~\ref{lb:lem:sep} and a strong limit give
$E(B)P(A)=0$.
If $P(\mathcal E)=I$, compact approximation of $A_0\cap\mathcal E$
for arbitrary disjoint compact $A_0,B_0\subset X$ gives
$E(B_0)P(A_0)=0$. For an open set $O$, exhaust $O$ by compact sets and use the two
orders of the preceding annihilation with the compact set $X\setminus O$.
Strong limits give
\[
 E(O)P(X\setminus O)=0,\qquad E(X\setminus O)P(O)=0.
\]
Since the two effects sum to $I$, these identities imply $E(O)=P(O)$.
Regularity of the scalar matrix-coefficient measures gives $E=P$.

For an atom $x$, every scalar representing measure with barycenter $x$
is $\delta_x$, since $x$ is extreme. This gives the unique extension
property for evaluation at $x$ with any multiplicity. Countable atomic
representations are direct sums of these.

If $F=\{b=0\}$ for an affine $b\ge0$ on $K$ and $P(F)=I$, then
$V^*\rho(b)V=\pi(b)=0$. Thus $Q(F)V=V$ and the dilation restricts to
$C(X\cap F)$. The identity $X\cap F=\ex F$ and compactness of this set
allow Scherer's planar theorem \cite[Theorem 3.8]{planar} to be applied; every affine function on $F$
extends to the ambient space.

Finally, a representation is a direct sum over a countable Borel
partition. If the compression to a summand $H_j$ is equal to its
representation on all of $C(X)$, the Stinespring isometry $V_j=V|_{H_j}$
satisfies $\rho(f)V_j=V_j\pi_j(f)$. This follows, for example, by
expanding the square of the difference using the equalities for
$f^*f,f,f^*$. Combining these intertwinings over the orthogonal sum
proves the assertion.
\end{proof}

\subsection{The dilation replacement and endpoint interpolation}

\begin{proposition}\label{prop:partner}
Every representation concentrated on $Z$ from Lemma~\ref{lem:measurable}
has the unique extension property.
\end{proposition}
\begin{proof}
By Lemma~\ref{lem:separable}, assume the source and target spaces
are separable.

Put $Y=X\setminus S$ and split the target space using its actual spectral
measure:
\[
 V_s=Q(S)V,\qquad V_y=Q(Y)V,\qquad A=V_y^*V_y.
\]
Write $\rho_s$ and $\rho_y$ for the restrictions of $\rho$ to these
reducing target spaces. These are contractions $V_s,V_y$, not generally
isometries.

\emph{Normal intertwining for the smooth target part.}
Let $P_n(D)=P(Z\cap n^{-1}(D))$ and
$Q_{s,n}(D)=Q(S\cap n^{-1}(D))$. If $x\in Z$ and $y\in S$
have distinct normals, their contact gap is strictly positive: equality
would make $n(x)$ a unit normal at the smooth point $y$.
Apply Lemma~\ref{lem:contact} to source and target sets with normals
in $D$ and $D^c$, respectively, and then reverse them. For every Borel
$D\subset S^2$ this gives
\[
                     Q_{s,n}(D)V_s=V_sP_n(D).
\]
Consequently $V_s^*V_s$ commutes with $P_n(D)$, and so does
$A=I-V_s^*V_s$.

\emph{A graph intertwiner for the nonsmooth target part.}
For $x\in Z$ and $y\in Y$, the equality
$\langle n(x),x-y\rangle=0$ holds precisely when $y=\tau(x)$.
Indeed, $X\cap F_{n(x)}=\{x,\tau(x)\}$ and $x\notin Y$.
For a Borel set $B\subset X$, apply Lemma~\ref{lem:contact} to
\[
 (Z\cap\tau^{-1}(B^c),\ Y\cap B)
 \quad\hbox{and}\quad
 (Z\cap\tau^{-1}(B),\ Y\setminus B).
\]
Both products have strictly positive contact gap. The resulting two
annihilations give the exact graph relation
\[
                   Q(B)V_y=V_yP(Z\cap\tau^{-1}(B)).
\]
Thus, defining a unital representation on the original source space by
\[
                \pi_\tau(f)=\int_Z f(\tau(x))\,dP(x),
\]
bounded Borel functional calculus gives
\[
                  \rho_y(f)V_y=V_y\pi_\tau(f),
 \qquad
 V_y^*\rho_y(f)V_y=A\pi_\tau(f)=\pi_\tau(f)A.
\]
Here bounded Borel functions on $Z$, such as $f\circ\tau$, may be
extended arbitrarily to $X\setminus Z$; the resulting operators are
independent of that extension because $P(Z)=I$.
In particular $A$ commutes with $\pi_\tau(C(X))$ as well as with the
normal algebra $P_n$. The formula remains valid when $\tau(x)$ is smooth:
the graph relation simply makes $V_y$ vanish on that source part.

\emph{Replacement of the target.}
On $Q(S)L\oplus H$ define
\[
 \widetilde\rho(f)=\rho_s(f)\oplus\pi_\tau(f),
 \qquad Wh=V_sh\oplus A^{1/2}h.
\]
Then $W^*W=V_s^*V_s+A=I$, and, for every $f\in C(X)$,
\[
 W^*\widetilde\rho(f)W
 =V_s^*\rho_s(f)V_s+A^{1/2}\pi_\tau(f)A^{1/2}
 =V_s^*\rho_s(f)V_s+V_y^*\rho_y(f)V_y
 =\Phi(f).
\]
Thus the entire UCP map is reproduced, not only its affine restriction.
Give the second target summand the \emph{source} normal label $n(x)$.
This is legitimate because both $\pi_\tau(C(X))$ and $A$ commute with
$P_n$. Explicitly,
\[
 (Q_{s,n}(D)\oplus P_n(D))W=WP_n(D).
\]
This is the missing normal intertwining. It does not require the normal
label to be a function of the nonsmooth target point $\tau(x)$.

\emph{Bounded interpolation on exposed segments.}
Use the Borel set $U$ and endpoints $a,b$ from
Lemma~\ref{lem:segment-endpoints}. Choose the first coordinate $j$
that varies on each segment, as in that lemma, and partition $U$
into countably many Borel pieces on each of which $j$ is fixed and
$b_j(u)-a_j(u)\ge1/m$ for some fixed positive integer $m$.
Such a disjoint partition is obtained by also choosing the least
$m$ for which this inequality holds. Since every $F_{n(x)}$, $x\in Z$,
is a segment, $P_n(U)=I$.

Write $\widetilde Q_n(D)=Q_{s,n}(D)\oplus P_n(D)$.
This spectral measure commutes with $\widetilde\rho(C(X))$ and
satisfies $\widetilde Q_n(D)W=WP_n(D)$. In particular
$\widetilde Q_n(U)W=W$, so its target complement may be discarded.
The countable partition then reduces both representations and splits
$W$ into isometries between their corresponding summands. Work on one
such piece for the next calculation. Simple-function approximation
of the projection intertwinings gives
\[
 c(\widetilde Q_n)W=Wc(P_n)
 \qquad\text{for every bounded Borel normal function }c.
\]
Here $c(P_n)=\int c(u)\,dP_n(u)$, and similarly on the target.

Fix $f\in C(X)$. On the chosen normal piece set
\[
 \beta(u)=\frac{f(b(u))-f(a(u))}{b_j(u)-a_j(u)},
 \qquad \alpha(u)=f(a(u))-\beta(u)a_j(u).
\]
Both coefficients are bounded and Borel: the denominator is at least
$1/m$, and $K$ is bounded. For either endpoint $t$ of $F_u$,
$f(t)=\alpha(u)+\beta(u)t_j$. This identity applies pointwise in all
three spectral realizations. The source uses $(u,t)=(n(x),x)$;
the smooth target uses $(n(y),y)$ with $y\in S$ and $n(y)\in U$;
the second target copy uses $(n(x),\tau(x))$. In each case
$t\in X\cap F_u=\{a(u),b(u)\}$.
Consequently bounded Borel functional calculus gives
\[
 \begin{split}
 \widetilde\rho(f)&=\alpha(\widetilde Q_n)
                  +\beta(\widetilde Q_n)\widetilde\rho(x_j),\\
 \pi(f)&=\alpha(P_n)+\beta(P_n)\pi(x_j).
 \end{split}
\]
The normal operators commute with the corresponding representations.
The adjoint of the intertwining identity for $\overline\beta$ gives
$W^*\beta(\widetilde Q_n)=\beta(P_n)W^*$, also when $f$ is complex-valued.
Using this intertwining and the affine compression identity yields
\[
 \begin{split}
 W^*\widetilde\rho(f)W
 &=\alpha(P_n)+\beta(P_n)W^*\widetilde\rho(x_j)W\\
 &=\alpha(P_n)+\beta(P_n)\pi(x_j)=\pi(f).
 \end{split}
\]
Combining the countably many normal pieces proves equality on the
whole source space. The replacement reproduces $\Phi$, so $\Phi=\pi$.
No disintegration is used: all identities are identities of bounded
Borel functions in the original spectral calculi. The lower bound
on the segment width is imposed separately on each piece; no uniform
lower bound on all segment lengths is assumed.
\end{proof}

\begin{proof}[Proof of Theorem~\ref{thm:three}]
After the separable reduction in Lemma~\ref{lem:separable}, partition
$X$ into $\mathcal E$, the countable set $C$, successively disjoint
parts of $X\cap F_j$, and the remainder $Z$. All sets are Borel by
Lemma~\ref{lem:measurable}. Lemma~\ref{lem:good} handles every part
except $Z$, and Proposition~\ref{prop:partner} handles $Z$.
The direct-sum assertion of Lemma~\ref{lem:good} concludes the proof.
\end{proof}

The same conclusion holds for every compact convex subset of $\R^3$
with closed extreme boundary: if
its affine dimension is at most two, apply the planar theorem after an
affine identification with a subset of $\R^2$.



\section{An explicit counterexample in affine dimension thirty-three}
\label{sec:construction}

\begin{theorem}\label{thm:main}
There is an explicitly specified compact convex body
\(K\subset\mathbb R^{33}\), with \(\aff K=\mathbb R^{33}\) and
compact extreme boundary \(X=\ex K\), for which the affine function
system \(A(K)|_X\) is not hyperrigid in \(C(X)\).
More precisely, there are a unital representation
\(\pi:C(X)\to B(H)\) and a UCP map \(\Psi:C(X)\to B(H)\) such that
\[
 \Psi|_{A(K)|_X}=\pi|_{A(K)|_X},\qquad \Psi(Q)\ne\pi(Q),
 \qquad Q(x,z)=|x|^2.
\]
The Hilbert space \(H\) is separable and infinite dimensional,
the spectral measure of \(\pi\) has no point atoms, and
\(\dim_{\mathbb C}A(K)=34\).
\end{theorem}

We use the Walsh dilation of Bilich--Scherer \cite[Section 3]{BS},
with a different three-point distribution. The operator calculation
is included in full. Its telescoping identity determines the height
of an encoded dilation point. A uniform scalar estimate places all
such points in a truncated epigraph with a strongly convex lower graph. An explicit
boundary kernel then gives the required UCP map.

\subsection{The Walsh compression and its height identity}

Fix \(R>1\) and \(1<b<R^2\), and put
\[
 a=b^{-1},\qquad
 \mu=\frac{b}{2R^2}(\delta_{-R}+\delta_R)
                  +\left(1-\frac b{R^2}\right)\delta_0.
\]
Thus \(\mu\) is symmetric and has second moment \(b\).
In the discrete alphabet \(\{-1,1,-R,0,R,*\}\), with a separate
star symbol, define
\[
 \begin{split}
 G&=\{-1,1\}^{\mathbb N},\qquad Y_0=\{-R,0,R\}^{\mathbb N},\\
 Y_k&=\{-1,1\}^{k-1}\times\{*\}\times
        \{-R,0,R\}^{\{k+1,k+2,\ldots\}}\quad(k\ge1),\\
 Y&=G\cup\bigcup_{k\ge0}Y_k.
 \end{split}
\]
The numerical coordinate \(x_n\) sends the star to zero.
Give \(G\) Haar probability \(m\). Give \(Y_k\) the product
probability \(\mu_k\) with uniform binary prefix, star at \(k\),
and \(\mu\)-distributed tail; on \(Y_0\), use
\(\mu_0=\mu^{\otimes\mathbb N}\).
The space \(Y\) is compact: unbounded star positions can accumulate
only in \(G\), and every fixed branch is compact.
Each coordinate is continuous, and \(Y_0\) is clopen.

Set
\[
 H=L^2(G,m),\qquad
 \mathcal L=\bigoplus_{k\ge0}L^2(Y_k,\mu_k).
\]
For a finite \(F\subset\mathbb N\), write
\(\chi_F=\prod_{n\in F}\varepsilon_n\) and
\(F_{>k}=F\cap\{k+1,k+2,\ldots\}\).
For \(k\notin F\), the functions
\[
                  \chi_F^{(k)}
        =b^{-|F_{>k}|/2}\prod_{n\in F}x_n
\]
form an orthonormal family in the \(k\)-th summand.
Define on the Walsh basis
\begin{equation}\label{eq:V}
 V\chi_F=a^{|F|/2}\chi_F^{(0)}
       +\sqrt{1-a}\sum_{k\in F}
            a^{|F_{>k}|/2}\chi_{F\setminus\{k\}}^{(k)}.
\end{equation}

\begin{lemma}\label{lem:walsh}
The map \(V:H\to\mathcal L\) is an isometry. For bounded Borel
functions on \(Y\), let \(\Phi(f)=V^*M_fV\). Then
\begin{equation}\label{eq:multilinear}
 \Phi\left(\prod_{n\in E}x_n\right)
                  =M_{\prod_{n\in E}\varepsilon_n}
              \qquad(E\subset\mathbb N\text{ finite}).
\end{equation}
Also \(\Phi(\mathbf1_{Y_0})\chi_F=a^{|F|}\chi_F\) and
\(V1=\mathbf1_{Y_0}\).
\end{lemma}
\begin{proof}
On branch \(k\), symmetry and independence make a matrix coefficient
of multiplication by \(\prod_{n\in E}x_n\) vanish unless the two
product labels differ exactly in \(E\) and \(k\notin E\).
In the nonzero case, its value is \(b^{|E_{>k}|/2}\): a tail
coordinate in the multiplier contributes its second moment
divided by \(\sqrt b\), while common coordinates contribute one
after normalization. On branch zero, the same calculation has no
star restriction. Consequently, for \(C_0=F\cap F'\),
\[
 \left\langle M_{\prod_E x_n}V\chi_F,V\chi_{F'}\right\rangle
 =\mathbf1_{\{E=F\triangle F'\}}
 \left(a^{|C_0|}+(1-a)\sum_{k\in C_0}a^{|(C_0)_{>k}|}\right).
\]
The parenthesis is the finite geometric sum
\(a^r+(1-a)\sum_{j=0}^{r-1}a^j=1\).
Taking \(E=\varnothing\) proves isometry on the dense Walsh span;
general \(E\) proves \eqref{eq:multilinear}.
The two final assertions follow by retaining the branch-zero term.
\end{proof}

Let \(\beta\) be an even nonnegative function on the numerical
digits, with \(\beta(\pm1)=0\). Define diagonal contractions
\[
                  A_n\chi_F=a^{|F_{>n}|}\chi_F\qquad(n\ge0).
\]
The parameters \(0<q<1\) and \(c>0\) will be chosen so that
\begin{equation}\label{eq:conditions}
 \begin{aligned}
 \int\bigl(r^2-\beta(r)-1\bigr)\,d\mu(r)
                                      &=c(q^2-1),\\
 a\bigl(R^2-\beta(R)-1\bigr)
       +(1-a)\bigl(-1-\beta(0)\bigr)&=c(q^2-a).
 \end{aligned}
\end{equation}

\begin{lemma}\label{lem:tel}
Under \eqref{eq:conditions},
\begin{equation}\label{eq:tel}
       \Phi\bigl(x_n^2-\beta(x_n)\bigr)-I
                       =c(q^2A_n-A_{n-1})\qquad(n\ge1).
\end{equation}
\end{lemma}
\begin{proof}
Parity makes the left side diagonal. Fix \(F\) and put
\(r=|F_{>n}|\). If \(n\notin F\), the squared weights of the
branches before \(n\) sum to \(a^r\); their expected multiplier
\(x_n^2-\beta(x_n)-1\) is the first left side of
\eqref{eq:conditions}. Later branches contribute zero, since their
\(n\)-th coordinate is binary.
If \(n\in F\), the earlier weights sum to \(a^{r+1}\).
Under the normalized \(x_n^2\)-weighted measure, the coordinate
has magnitude \(R\), so its multiplier is
\(R^2-\beta(R)-1\). The star branch has weight
\((1-a)a^r\) and multiplier \(-1-\beta(0)\).
The two diagonal coefficients are therefore \(a^r\) times the
respective left sides of \eqref{eq:conditions}. These are exactly
the coefficients of the right side of \eqref{eq:tel}.
\end{proof}

Let \(D\) be a positive integer, let
\(j(n)=1+((n-1)\bmod D)\), and define
\begin{equation}\label{eq:encoding}
 \begin{split}
 w_n&=q^{n-1}e_{j(n)},\qquad
 t(y)=\sum_{n\ge1}w_nx_n(y),\qquad s=t|_G,\\
 U(y)&=\sum_{n\ge1}q^{2n-2}\beta(x_n(y)),\\
 z_Y(y)&=|t(y)|^2-U(y)+c\mathbf1_{Y_0}(y).
 \end{split}
\end{equation}
These series converge uniformly and absolutely. All the displayed
functions are continuous, and \(U|_G=0\).
The operator calculation already gives
\begin{equation}\label{eq:height}
                  \Phi(t_j)=M_{s_j},\qquad
                  \Phi(z_Y)=M_{|s|^2}.
\end{equation}
Indeed, all mixed terms in the square expansion of \(|t|^2\)
compress exactly by \eqref{eq:multilinear}. Termwise compression
is valid because \(\sum_n|w_n|<\infty\).
For the diagonal terms, multiplying \eqref{eq:tel} by
\(q^{2n-2}\) gives the norm-convergent telescope
\[
 \sum_{n=1}^Nq^{2n-2}
       \left[\Phi\bigl(x_n^2-\beta(x_n)\bigr)-I\right]
       =c(q^{2N}A_N-A_0)\longrightarrow-cA_0.
\]
This cancels \(c\Phi(\mathbf1_{Y_0})=cA_0\), proving
\eqref{eq:height}. No orthogonality or cancellation of distinct
vectors \(w_n\) is used in this calculation.

\subsection{Numerical choices and uniform containment}

The numerical data are collected below; the other
quantities retain their definitions above.
\begin{center}
\begin{tabular}{c@{\quad}c@{\qquad\qquad}c@{\quad}c}
\toprule
Quantity & Value & Quantity & Value\\ \midrule
\(D\) & \(32\) & \(R\) & \(12\)\\
\(b\) & \(21/2\) & \(q^2\) & \(45026/72405\)\\
\(\beta(0)\) & \(197/200\) & \(\beta(R)\) & \(2417/20\)\\
\(c\) & \(14481/24320\) & Distance coefficient & \(1-10^{-5}\)\\
\bottomrule
\end{tabular}
\end{center}
In particular \(\beta(\pm1)=0\) completes the specification of
\(\beta\). The two scalar identities \eqref{eq:conditions} are
verified exactly: their left sides are respectively
\[
                 -\frac{1441}{6400}=c(q^2-1),\qquad
                   \frac{439}{1400}=c(q^2-a).
\]

For the geometric estimate write \(\rho=q^D\) and
\[
 C_\rho=\left\{\sum_{\ell\ge0}\rho^\ell\varepsilon_\ell:
                              \varepsilon_\ell\in\{-1,1\}\right\}.
\]

\begin{lemma}[Uniform digit estimate]\label{lem:contain}
For every sequence \(d_\ell\in\{-R,-1,0,1,R\}\),
\begin{equation}\label{eq:scalarcontain}
 (1-10^{-5})\dist\left(\sum_{\ell\ge0}\rho^\ell d_\ell,C_\rho\right)^2
             \ge\sum_{\ell\ge0}\rho^{2\ell}\beta(d_\ell).
\end{equation}
\end{lemma}
\begin{proof}
Both sides vanish for an entirely binary sequence. Otherwise let
\(h\) be its first nonbinary position and put
\(\delta=\dist(d_h,\{-1,1\})\in\{1,R-1\}\).
The largest possible difference between a numerical digit and a
binary digit is \(R+1\). Put
\[
 T=\frac{(R+1)\rho}{1-\rho},\qquad
 E=\frac{\beta(R)\rho^2}{1-\rho^2}.
\]
Against any signed expansion, let \(m\le h\) be the first
mismatch. Its leading term dominates the remaining tail and gives
a distance at least
\(\rho^m(|d_m-\varepsilon_m|-T)\).
If \(m=h\), this is at least \(\rho^h(\delta-T)\).
If \(m<h\), it is at least \(\rho^h(R-1-T)\), provided
\(2-T\ge\rho(R-1-T)\). Thus the distance to \(C_\rho\)
is at least \(\rho^h(\delta-T)\).
The right side of \eqref{eq:scalarcontain} is at most
\(\rho^{2h}(\beta(d_h)+E)\), since \(\beta(R)\) is the
largest digit correction.

Here all the required estimates admit rational verification.
The integer inequality
\[
                  1990\cdot45026^{16}<72405^{16}
\]
gives \(\rho<1/1990\), hence
\[
 T<\frac{13}{1989}<1,\qquad
 E<\frac{2417}{20(1990^2-1)}.
\]
Also \(2-13/1989>11/1990\), which implies the earlier-mismatch
comparison \(2-T>\rho(R-1-T)\). Finally, clearing positive
denominators verifies
\[
 \begin{aligned}
 (1-10^{-5})\left(1-\frac{13}{1989}\right)^2
 -\frac{197}{200}-\frac{2417}{20(1990^2-1)}
                                      &>\frac{19}{10000},\\
 (1-10^{-5})\left(11-\frac{13}{1989}\right)^2
 -\frac{2417}{20}-\frac{2417}{20(1990^2-1)}
                                      &>\frac{49}{10000}.
 \end{aligned}
\]
These prove \((1-10^{-5})(\delta-T)^2>\beta(d_h)+E\)
in both cases. The estimates bound the entire infinite tail.
\end{proof}

\subsection{The convex body and its whole extreme boundary}

The represented horizontal set is the Cartesian product
\begin{equation}\label{eq:product}
                 C=s(G)=\prod_{j=1}^Dq^{j-1}C_\rho.
\end{equation}
Define
\begin{equation}\label{eq:g}
 \begin{split}
 g(x)&=|x|^2-(1-10^{-5})\dist(x,C)^2\\
 &=10^{-5}|x|^2+(1-10^{-5})
                  \sup_{u\in C}(2\langle x,u\rangle-|u|^2).
 \end{split}
\end{equation}
The second formula proves strong convexity: it is a positive
quadratic plus a supremum of affine functions.
With
\[
 B=[-R-1,R+1]^D,\qquad M=(R+1)^2D+c+1,
\]
let
\begin{equation}\label{eq:K}
              K=\{(x,z):x\in B,\quad g(x)\le z\le M\}.
\end{equation}
This set is compact and convex. Since
\(g(x)\le|x|^2\le(R+1)^2D<M\) throughout \(B\),
it has nonempty interior in \(\mathbb R^{D+1}\).

\begin{lemma}\label{lem:extreme}
The entire extreme boundary is the compact set
\begin{equation}\label{eq:X}
 X=\{(x,g(x)):x\in B\}
           \cup\{((R+1)\sigma,M):\sigma\in\{-1,1\}^D\}.
\end{equation}
Moreover the continuous maps
\[
 \iota_G(\varepsilon)=(s(\varepsilon),|s(\varepsilon)|^2),
 \qquad \iota_Y(y)=(t(y),z_Y(y))
\]
take their values in \(X\) and \(K\), respectively.
\end{lemma}
\begin{proof}
In a convex decomposition of a lower-graph point, equality in
the convexity inequality forces the horizontal coordinates to
coincide by strict convexity, and then forces the heights to
coincide. Every lower-graph point is therefore extreme.
An intermediate-height point is a nontrivial vertical convex
combination. The top face is a box, whose only extreme points
are its vertices. This proves \eqref{eq:X}, including the absence
of any other extreme points. Compactness follows from continuity
of \(g\) and finiteness of the vertex set.

For the containment assertion,
\( |t_j(y)|\le Rq^{j-1}/(1-\rho)<R+1\).
Also \(U\ge0\), so \(z_Y\le|t|^2+c<M\).
Apply Lemma~\ref{lem:contain} to each residue class of digits
and multiply by \(q^{2(j-1)}\). Squared distance to the product
\eqref{eq:product} is the sum of the squared coordinate distances,
so
\[
                 (1-10^{-5})\dist(t(y),C)^2\ge U(y).
\]
Consequently \(z_Y(y)\ge g(t(y))\), proving \(\iota_Y(Y)\subset K\).
Finally \(g=|\cdot|^2\) on \(C\), so \(\iota_G(G)\subset X\).
\end{proof}

\subsection{The boundary kernel and a positive witness}

For \(w=(x,z)\in K\), put
\[
 r(w)=\frac{z-g(x)}{M-g(x)},\qquad
 p_\sigma(x)=\prod_{j=1}^D\frac{1+\sigma_jx_j/(R+1)}2.
\]
The following probability measure on \(X\) has barycenter \(w\):
\begin{equation}\label{eq:kernel}
 \nu_w=(1-r(w))\delta_{(x,g(x))}
       +r(w)\sum_{\sigma\in\{-1,1\}^D}
                p_\sigma(x)\delta_{((R+1)\sigma,M)}.
\end{equation}
Indeed all weights are nonnegative and sum to one, and
\(\sum_\sigma p_\sigma(x)(R+1)\sigma=x\).
The denominator \(M-g(x)\) is positive on the whole box.
For every \(f\in C(X)\), the function
\(w\mapsto\int f\,d\nu_w\) is continuous.

Define
\begin{equation}\label{eq:maps}
 \pi(f)=M_{f\circ\iota_G},\qquad
 \widehat f(y)=\int_X f\,d\nu_{\iota_Y(y)},\qquad
 \Psi(f)=V^*M_{\widehat f}V.
\end{equation}
The first map is a unital representation. Pointwise integration
against probability measures preserves positivity at every matrix
level; multiplication and isometric compression therefore make
\(\Psi\) UCP. This formula uses the explicit finite kernel
\eqref{eq:kernel}, with no choice of representing measures.

It also gives a representation dilation directly. Index the terms
of \eqref{eq:kernel} by \(J=\{0\}\cup\{-1,1\}^D\), writing
their locations and weights at \(\iota_Y(y)\) as
\(F_\ell(y)\in X\) and \(\omega_\ell(y)\ge0\).
On
\(\widetilde{\mathcal L}
 =\bigoplus_{k\ge0}\bigoplus_{\ell\in J}L^2(Y_k,\mu_k)\),
let \(\rho_X(f)\) multiply component \((k,\ell)\) by
\(f(F_\ell(y))\). The map
\[
              (W\xi)_{k,\ell}(y)
                    =\sqrt{\omega_\ell(y)}\,\xi_k(y)
\]
is an isometry, and
\(\Psi(f)=(WV)^*\rho_X(f)(WV)\).

\begin{proof}[Completion of the proof of Theorem~\ref{thm:main}]
For an affine function, its integral against \(\nu_w\) is its
value at \(w\). Thus \eqref{eq:height} gives
\[
 \Psi(x_j|_X)=M_{s_j}=\pi(x_j|_X),\qquad
 \Psi(z|_X)=M_{|s|^2}=\pi(z|_X).
\]
Together with constants, these coordinates span \(A(K)|_X\).
They are linearly independent because \(K\) has full affine
dimension \(D+1=33\), and restriction to \(X\) is injective
because \(K=\operatorname{conv}X\). Hence the function system
has complex dimension \(D+2=34\).

For \(Q(x,z)=|x|^2\), scalar convexity and the barycenter
property give \(\widehat Q(y)\ge|t(y)|^2\).
The constant Walsh vector has norm one and
\(V1=\mathbf1_{Y_0}\). Independence and zero means remove all
mixed terms from the following scalar integrals:
\[
 \begin{split}
 \langle\Psi(Q)1,1\rangle
       &\ge\int_{Y_0}|t|^2\,d\mu_0=\frac b{1-q^2},\\
 \langle\pi(Q)1,1\rangle
       &=\int_G|s|^2\,dm=\frac1{1-q^2}.
 \end{split}
\]
Consequently
\[
              \langle[\Psi(Q)-\pi(Q)]1,1\rangle
                         \ge\frac{b-1}{1-q^2}>0.
\]
The representation therefore fails the unique extension property.

Finally \(q^D=\rho<1/2\) makes each signed expansion in
\(C_\rho\) injective: the first differing digit exceeds the
entire remaining difference tail. The product encoding \(s\),
and hence \(\iota_G\), is injective. A point inverse image under
\(\iota_G\) is therefore empty or a singleton; Haar measure on
\(G\) assigns it measure zero. This proves that the spectral
measure of \(\pi\) has no point atoms. The countably infinite
Walsh basis proves that \(H\) is separable and infinite dimensional.
\end{proof}

\section{An obstruction for stationary geometric Walsh encodings}
\label{enc:sec}

The construction above separates the operator identities from the
geometric containment problem. We now show that its dimension
cannot be reduced to four by changing only the fixed three-point
distribution, the geometric decay rate, or the directions of the
encoding vectors. The first obstruction concerns the quadratic
distance envelope. The next section extends it, in affine dimension
four, to every finite strictly convex lower graph.

\subsection{A signing estimate}

\begin{lemma}[Signed series in finite dimension]\label{enc:signed}
Let $w_n\in\mathbb R^D$ satisfy $\sum_n|w_n|<\infty$, and let
$L_D$ be the sum of the $D$ largest numbers among
$\{|w_n|^2:n\ge1\}$, padding with zeros if necessary.
There are signs $\varepsilon_n\in\{-1,1\}$ such that
\[
 \left|\sum_{n\ge1}\varepsilon_nw_n\right|^2\le L_D.
\]
\end{lemma}
\begin{proof}
For a finite sequence $w_1,\ldots,w_N$, the polytope
\[
 \left\{a\in[-1,1]^N:\sum_{n=1}^N a_nw_n=0\right\}
\]
is nonempty. At an extreme point $a$, at most $D$ coordinates
lie strictly between $-1$ and $1$: otherwise the corresponding
vectors are linearly dependent and admit a nonzero feasible
perturbation in both directions. Independently round each of
these coordinates to a sign with mean $a_n$, leaving the other
coordinates fixed. Since the centered rounding errors are
independent,
\[
 \mathbb E\left|\sum_{n=1}^N\varepsilon_nw_n\right|^2
 =\sum_{|a_n|<1}(1-a_n^2)|w_n|^2\le L_D.
\]
Some realization therefore has the required bound. Apply this
argument to finite prefixes and extend their signings arbitrarily.
The remaining tail has norm at most $\sum_{n>N}|w_n|\to0$.
Compactness of the sign space and uniform convergence of the
signed series give the assertion for the infinite sequence.
\end{proof}

\subsection{Eliminating all fixed three-point parameters}

In this section, $\gamma$ denotes the squared geometric ratio;
thus $\gamma=q^2$ in the construction above. The numerical tail
law is supported on $\{0,\pm R\}$, with variance $b$ and masses
$1-b/R^2,b/(2R^2),b/(2R^2)$. As before, $G$ is the binary source,
$Y_0$ has only tail coordinates, and $Y_k$ has a binary prefix,
a zero-valued star at $k$, and tail coordinates thereafter.

\begin{theorem}[No stationary geometric distance encoding in
horizontal dimension at most four]\label{enc:universal}
Let
\[
 R>1,\qquad 1<b<R^2,\qquad c>0,\qquad 0<\gamma<1,
 \qquad u,v\ge0.
\]
Suppose the exact height equations hold:
\begin{equation}\label{enc:eq:telescoping}
\begin{aligned}
 \frac b{R^2}(R^2-1-v)
  +\left(1-\frac b{R^2}\right)(-1-u)&=c(\gamma-1),\\
 \frac1b(R^2-1-v)
  +\left(1-\frac1b\right)(-1-u)&=c(\gamma-1/b).
\end{aligned}
\end{equation}
Let $w_n\in\mathbb R^D$ have arbitrary directions and satisfy
$|w_n|^2=\gamma^{n-1}$. Define
\[
 t(y)=\sum_{n\ge1}w_nx_n(y),\qquad C=t(G),\qquad
 U(y)=\sum_{n\ge1}\gamma^{n-1}\beta(x_n(y)),
\]
where $\beta(\pm1)=0$, $\beta(0)=u$, and $\beta(\pm R)=v$.
If $D\le4$, there is no $0<\lambda\le1$ for which
\begin{equation}\label{enc:eq:containment}
 \lambda\dist(t(y),C)^2-U(y)+c\mathbf1_{Y_0}(y)\ge0
 \qquad(y\in Y).
\end{equation}
\end{theorem}
\begin{proof}
Two consequences of Lemma~\ref{enc:signed} suffice. At a star in
the first position followed by zero digits, the encoded point is
zero and $U=u/(1-\gamma)$. Signing the source series bounds its
distance to $C$. On $Y_0$, take digits $R\varepsilon_n$ and compare
with the source point having digits $\varepsilon_n$; the difference
is $(R-1)\sum_n\varepsilon_nw_n$. Thus
\eqref{enc:eq:containment} implies, after discarding $\lambda\le1$,
\begin{equation}\label{enc:eq:necessary}
 u\le1-\gamma^D,\qquad
 v-c(1-\gamma)\le(R-1)^2(1-\gamma^D).
\end{equation}
For $D\le4$ these inequalities remain necessary with $D$ replaced
by four.

Put $x=\gamma^4$, $h=1+u$, and $\Delta=R^2+u-v$.
The height equations give
\begin{equation}\label{enc:eq:elimination}
 c=\frac{\Delta(R^2-b^2)}{R^2(b-1)},\qquad
 c(1-\gamma)=h-\frac b{R^2}\Delta.
\end{equation}
Using the second equality in the second inequality of
\eqref{enc:eq:necessary} yields
\begin{equation}\label{enc:eq:positive-delta}
 \left(1-\frac b{R^2}\right)\Delta
 \ge2(R-1)+(R-1)^2x>0.
\end{equation}
Hence $\Delta>0$, and the first equality of
\eqref{enc:eq:elimination}, together with $c>0$, forces $b<R$.
In particular all subsequent denominators are positive.
Eliminating $c$ also gives
\[
 \frac h\Delta
 =\frac{R^2-b-\gamma(R^2-b^2)}{R^2(b-1)}>0.
\]
Multiply \eqref{enc:eq:positive-delta} by $h/\Delta$ and use
$h\le2-x$, from the first inequality in
\eqref{enc:eq:necessary}. After cancellation,
\begin{equation}\label{enc:eq:reduced}
 (R-1)(2+(R-1)x)
 \left(1-\gamma\frac{R^2-b^2}{R^2-b}\right)
 \le(b-1)(2-x).
\end{equation}

Set $r=R-1>0$, $\alpha=(b-1)/(R-1)\in(0,1)$, and
$A=1-\alpha\in(0,1)$. Clearing the positive denominators in
\eqref{enc:eq:reduced} says that $p(r)\ge0$, where
\[
 p(r)=\gamma(1-\alpha)(2+(1+\alpha)r)(2+rx)
       -(r+2-\alpha)(2(1-\alpha)+(r+\alpha)x).
\]
Its quadratic, linear, and constant coefficients are respectively
\begin{equation}\label{enc:eq:coefficients}
\begin{aligned}
 &x\bigl(\gamma(1-\alpha^2)-1\bigr),\\
 &2\bigl[-\gamma A^2+(2\gamma-1+\gamma x)A-x\bigr],\\
 &-(2-x)A^2+2(2\gamma-1)A-x.
\end{aligned}
\end{equation}
We prove that all three are strictly negative.
The first is immediate. For the second put
$N=2\gamma-1+\gamma^5$. If $N\le0$, negativity is immediate;
if $N>0$, half the coefficient is at most
$N^2/(4\gamma)-\gamma^4$. We claim that
\begin{equation}\label{enc:eq:linear-bound}
 2\gamma-1+\gamma^5<2\gamma^{5/2}\qquad(0<\gamma<1).
\end{equation}
Indeed, setting $t=\sqrt\gamma$, the difference factors as
\[
 t^{10}-2t^5+2t^2-1=(t-1)Q(t),
\]
where
\[
 Q(t)=t^9+t^8+t^7+t^6+t^5-t^4-t^3-t^2+t+1>0
 \qquad(0<t<1).
\]
For $t\le2/3$, this follows from
$t^2+t^3+t^4\le76/81<1$. For $t\ge2/3$, use
\[
 Q(t)=(1-t^2)(1+t)
       +t^4(t+t^2+t^3+t^4+t^5-1),
\]
whose first term is positive and whose second is nonnegative
because $t+t^2\ge10/9$. This proves
\eqref{enc:eq:linear-bound} and negativity of the linear coefficient.

For the constant coefficient in \eqref{enc:eq:coefficients},
$\gamma\le1/2$ is immediate. If $\gamma>1/2$, its maximum over
real $A$ is
\[
 \frac{(2\gamma-1)^2}{2-\gamma^4}-\gamma^4<0.
\]
To verify the strict inequality, factor its numerator as
\[
\begin{split}
 (2\gamma-1)^2-2\gamma^4+\gamma^8
 ={}&(\gamma-1)\bigl(\gamma^7+\gamma^6+\gamma^5+\gamma^4\\
 &\hspace{43mm}-\gamma^3-\gamma^2+3\gamma-1\bigr).
\end{split}
\]
The parenthesis is positive, since
$3\gamma-1-\gamma^2-\gamma^3
=(1-\gamma)(\gamma^2+2\gamma-1)\ge0$ on this interval and the
remaining powers are positive. All coefficients of $p$ are
therefore negative, so $p(r)<0$ for $r>0$, a contradiction.
\end{proof}

Theorem~\ref{enc:universal} permits arbitrary, nonperiodic directions
and every fixed symmetric three-point distribution with nonnegative
potential satisfying the stated height equations. It excludes the distance-envelope realization
in affine dimensions four and five, but is not a universal lower
bound for commutative counterexamples. The remaining possibility
of changing the lower graph requires an operator argument.

\section{Strictly convex graphs in affine dimension four}
\label{prox:sec}

The quadratic envelope is forced by an injective affine defect
when the horizontal dimension is at most three. The proof requires
no smoothness of the lower graph. Away from its quadratic
tangencies, supporting affine functions admit a metric taking
values on a positive distance level. Such levels have the critical
covering bound needed by Theorem~\ref{bar:thm}.

\subsection{Distance levels and supporting affine functions}

\begin{lemma}[Covering a positive distance level]\label{prox:lem:level}
Let $C\subset\mathbb R^D$ be nonempty and compact, and fix
$\varepsilon>0$. The compact set
\[
 S_\varepsilon=\{z:\dist(z,C)=\varepsilon\}
\]
has covers by $O(r^{-(D-1)})$ sets of diameter at most $r$ as
$r\downarrow0$.
\end{lemma}
\begin{proof}
For $z\in S_\varepsilon$, choose a nearest point $c\in C$ and put
$v=(z-c)/\varepsilon$. If $z'\in S_\varepsilon$ and $d=z'-z$,
then $|z'-c|\ge\varepsilon$. Interchanging the two points gives,
for any nearest-point unit vector $v'$ at $z'$,
\begin{equation}\label{prox:eq:normal-bounds}
 v\cdot d\ge-\frac{|d|^2}{2\varepsilon},\qquad
 v'\cdot d\le\frac{|d|^2}{2\varepsilon}.
\end{equation}
Cover $S_\varepsilon$ by finitely many closed spatial sets of
diameter at most $\varepsilon/4$, and the unit sphere by finitely
many closed balls of radius $1/8$ with unit centers $e_j$.
Partition into overlapping patches by requiring membership in
one spatial set and the existence of a nearest-point vector
within $1/8$ of one $e_j$. These patches are compact, since they
are projections of compact sets of nearest-point pairs.

For two points in the same patch, choose witnessing vectors in
\eqref{prox:eq:normal-bounds}. Then
\[
 |e_j\cdot d|\le\frac18|d|+
           \frac{|d|^2}{2\varepsilon}\le\frac14|d|.
\]
Orthogonal projection onto $e_j^\perp$ is therefore injective on
the patch and has a Lipschitz inverse. A grid in its bounded
projected subset of $\mathbb R^{D-1}$ gives the stated covering
bound. Sum over the finitely many patches. For $D=1$, each
patch contains at most one point. No continuity or uniqueness
of nearest-point vectors is needed.
\end{proof}

For a finite convex function $g$ on an open convex set $\Omega$, we
write
\[
 \partial g(s)=\{p:\ g(t)\ge g(s)+\langle p,t-s\rangle
                         \text{ for all }t\in\Omega\}
\]
for its convex subdifferential. The next lemma uses subgradients that
differ from the gradient $2s$ of the quadratic source graph.

\begin{lemma}[A metric for subgradient barriers]\label{prox:lem:metric}
Let $g$ be finite and convex on an open convex set
$\Omega\subset\mathbb R^D$, and let $C\subset\Omega$ be compact
with $g(s)=|s|^2$ on $C$. Suppose $F\subset C$ is Borel and
$p_s\in\partial g(s)$ is a Borel selection satisfying
$0<\delta\le|p_s-2s|\le M<\infty$ on $F$.
Set
\[
 v_s=\frac{p_s-2s}{|p_s-2s|},\qquad
 \varepsilon=\delta/16,\qquad \theta_s=s+\varepsilon v_s.
\]
Then $\theta:F\to S_\varepsilon$ is Borel, with the distance
level taken relative to the whole set $C$, and for a finite
constant $L$ depending only on $\delta,M$,
\begin{equation}\label{prox:eq:barrier-bound}
 0\le|t|^2-|s|^2-p_s\cdot(t-s)
       \le L|\theta_t-\theta_s|^2\qquad(s,t\in F).
\end{equation}
\end{lemma}
\begin{proof}
The subgradient inequality and $g=|\cdot|^2$ on $C$ give
\[
 (p_s-2s)\cdot(t-s)\le|t-s|^2\qquad(s\in F,\ t\in C).
\]
Consequently
\[
 |\theta_s-t|^2\ge\varepsilon^2+
       (1-2\varepsilon/\delta)|t-s|^2.
\]
Thus $s$ is the unique nearest point in $C$ to $\theta_s$, and
$\theta_s\in S_\varepsilon$.

For $s,t\in F$ write $d=t-s$ and $w=v_t-v_s$.
Applying the subgradient inequality in both orders gives
\[
 v_s\cdot d\le|d|^2/\delta,\qquad
 v_t\cdot d\ge-|d|^2/\delta,
\]
so $w\cdot d\ge-2|d|^2/\delta$ and
\begin{equation}\label{prox:eq:metric-control}
 |\theta_t-\theta_s|^2
 =|d+\varepsilon w|^2
 \ge\tfrac34|d|^2+\varepsilon^2|w|^2.
\end{equation}
The barrier in \eqref{prox:eq:barrier-bound} is nonnegative by
convexity. Writing $\alpha_s=|p_s-2s|\le M$, its upper bound is
\[
 |d|^2-\alpha_s v_s\cdot d
 \le(1+M/\delta)|d|^2+M|w|\,|d|.
\]
Young's inequality and \eqref{prox:eq:metric-control} bound this
by $L|\theta_t-\theta_s|^2$ for a finite $L$.
\end{proof}

\subsection{An injective defect forces the quadratic envelope}

\begin{theorem}[An injective defect forces the quadratic envelope]
\label{prox:thm:envelope}
Let $1\le D\le3$, and let $g$ be finite and strictly convex on an
open convex set $\Omega\subset\mathbb R^D$.
Let $K\subset\Omega\times\mathbb R$ be compact and convex, with
$X=\ex K$ compact and $K\subset\{(t,z):z\ge g(t)\}$.
Suppose $C\subset\Omega$ is nonempty and compact, and
\[
 g(s)=|s|^2,\qquad x_s=(s,|s|^2)\in K\qquad(s\in C).
\]
The graph points form a set $E=\{x_s:s\in C\}\subset X$.
Let $\pi:C(X)\to B(H)$ be a unital representation concentrated
on $E$ and having full spectral support there.
Suppose a UCP map $\Psi:C(X)\to B(H)$ agrees with $\pi$ on
affine functions, and for some real affine $h$ the operator
\[
 \mathcal D_h=\Psi(h^2)-\pi(h)^2
\]
is injective. Then $\pi$ is concentrated on points where
$\partial g(s)=\{2s\}$, equivalently where $g$ is differentiable
with $\nabla g(s)=2s$. Moreover,
\begin{equation}\label{prox:eq:envelope}
 2s\in\partial g(s)\quad(s\in C),\qquad
 g(t)\ge|t|^2-\dist(t,C)^2\quad(t\in\Omega).
\end{equation}
\end{theorem}
\begin{proof}
Subdifferentials of a finite convex function on an open convex
domain are nonempty and compact, locally bounded, and have closed
graph. Local boundedness follows by evaluating the supporting
inequality a fixed small distance in every direction inside a
larger compact ball. Hence the subgradient graph over $C$ is
compact, and $|p-2s|\le M$ there for some finite $M$.

For $m\ge1$ let
\[
 \Gamma_m=\{(s,p):s\in C,\ p\in\partial g(s),\
                    |p-2s|\ge1/m\},\qquad
 F_m=\operatorname{pr}_C\Gamma_m.
\]
These sets are compact. The increasing sets $F_m$ exhaust
$\{s\in C:\partial g(s)\ne\{2s\}\}$.
On each nonempty $F_m$, Lemma~\ref{lb:lem:selection} gives a
Borel selector with $1/m\le|p_s-2s|\le M$.
For the graph $E_m=\{x_s:s\in F_m\}$ define
\[
 b_{x_s}(t,z)=z-|s|^2-p_s\cdot(t-s).
\]
The coefficients are Borel, and $b_{x_s}$ is nonnegative on $K$
and vanishes there only at $x_s$. Indeed strict convexity makes
the subgradient inequality strict at every $t\ne s$; at $t=s$,
equality requires $z=g(s)$. This also shows that all graph points
under consideration are exposed extreme points of $K$.

Lemma~\ref{prox:lem:metric}, with $\delta=1/m$, bounds these
barriers on $E_m$ by a constant times squared distance after a
Borel map into $S_{1/(16m)}$. By Lemma~\ref{prox:lem:level}, this
compact level has $O(r^{-(D-1)})$ covers and hence, since $D\le3$,
$O(r^{-2})$ covers. Lemma~\ref{lb:lem:energy} gives infinite
inverse-square energy for every nonzero finite measure pushed
onto it. Theorem~\ref{bar:thm} therefore gives UEP for every
representation concentrated on $E_m$.

Write $P$ for the spectral measure of $\pi$ and $P_m=P(E_m)$.
Restrict $\pi$ and compress $\Psi$ to each nonzero $P_mH$.
These maps still agree on the affine functions, so the preceding
UEP assertion makes them equal. Since $P_m$ commutes with
$\pi(h)$,
\[
 P_m\mathcal D_hP_m=0.
\]
The operator $\mathcal D_h$ is positive by the Schwarz inequality.
Thus $\mathcal D_h^{1/2}P_m=0$, and injectivity forces $P_m=0$.
The Borel set
\[
 E_{\mathrm{bad}}=\bigcup_mE_m
 =\{x_s:s\in C,\ \partial g(s)\ne\{2s\}\}
\]
has spectral projection zero. A finite convex function is
differentiable precisely where its subdifferential is a singleton,
which proves the concentration assertion.

Full spectral support makes the remaining source points dense
in $C$. The closed subgradient graph then gives
$2s\in\partial g(s)$ for every $s\in C$. Finally,
\[
 g(t)\ge\sup_{s\in C}
       \bigl(g(s)+2s\cdot(t-s)\bigr)
       =|t|^2-\dist(t,C)^2,
\]
proving \eqref{prox:eq:envelope}.
\end{proof}

\begin{corollary}[UEP away from quadratic tangencies]
\label{prox:cor:source}
Under the geometric hypotheses of Theorem~\ref{prox:thm:envelope},
every representation concentrated on
$E_{\mathrm{bad}}=\{x_s:s\in C,\ \partial g(s)\ne\{2s\}\}$
has UEP. More generally, for any representation $\pi$ concentrated
on $E$ and any affine-agreeing UCP map $\Psi$,
\[
 \bigl(\Psi(h^2)-\pi(h)^2\bigr)P(E_{\mathrm{bad}})=0
\]
for every real affine $h$.
\end{corollary}
\begin{proof}
Use the increasing compact sets $E_m$ in the preceding proof,
each of which supports only representations with UEP.
If $P(E_{\mathrm{bad}})=I$, then $P(E_m)\uparrow I$ and UEP on
the compressed spaces gives
\[
 P(E_m)\Psi(f)P(E_m)=P(E_m)\pi(f)P(E_m)\qquad(f\in C(X)).
\]
Taking strong limits proves $\Psi=\pi$.
Without that concentration assumption, compressed equality still
gives $P(E_m)\mathcal D_hP(E_m)=0$. Positivity yields
$\mathcal D_hP(E_m)=0$, and the strong limit gives the result.
\end{proof}

\subsection{Application to the Walsh construction}

\begin{lemma}[Injective horizontal defects]\label{reg:lem:injective}
For the stationary three-point Walsh compression, assume
$1<b<R^2$ and write $\Phi(f)=V^*M_fV$.
Each operator $D_n=\Phi(x_n^2)-I$ is positive and injective.
If $t_j=\sum_n(w_n)_jx_n$ is a nonzero linear coordinate with
$\sum_nR|w_n|<\infty$ and $s_j=t_j|_G$, then
\begin{equation}\label{reg:eq:defect-sum}
 \Phi(t_j^2)-M_{s_j}^2=\sum_n(w_n)_j^2D_n
\end{equation}
is also positive and injective. A positive barycentric-kernel
realization with horizontal coordinate $t_j$ has an actual
coordinate defect at least as large as this operator.
\end{lemma}
\begin{proof}
Put $a=1/b$. On the Walsh vector $\chi_F$, let $r=|F_{>n}|$.
If $n\notin F$, the total squared weight of branches before
$n$ is $a^r$, and their tail second moment is $b$; later branches
have binary value squared equal to one. If $n\in F$, the total
weight before $n$ is $a^{r+1}$, their squarefree-product-weighted
second moment is $R^2$, and branch $n$ has weight
$(1-a)a^r$ and value zero. Consequently
\[
 D_n\chi_F=a^r
 \begin{cases}
 b-1,&n\notin F,\\
 R^2/b-1,&n\in F.
 \end{cases}\chi_F.
\]
Every diagonal entry is strictly positive, so $D_n$ is injective;
no uniform positive lower bound is asserted.

Uniform absolute convergence permits expansion of $t_j^2$.
All mixed terms compress exactly by
\eqref{eq:multilinear}, leaving the norm-convergent positive sum
\eqref{reg:eq:defect-sum}. One nonzero coefficient bounds this
sum below by a positive scalar multiple of an injective $D_n$.
Finally, scalar Jensen for a probability measure with barycenter
$(t(y),z(y))$ gives $\int\tau_j^2\,d\nu_y\ge t_j(y)^2$.
Applying the positive compression shows
$\Psi(\tau_j^2)\ge\Phi(t_j^2)$, while affine agreement gives
$\pi(\tau_j)=M_{s_j}$. This proves the last assertion.
\end{proof}

\begin{corollary}[No strictly convex graph in affine dimension four]
\label{prox:cor:walsh}
Take the stationary geometric data of
Theorem~\ref{enc:universal}. Suppose a compact convex body $K$
with compact extreme boundary contains every encoded point
\[
 (t(y),z_Y(y)),\qquad
 z_Y(y)=|t(y)|^2-U(y)+c\mathbf1_{Y_0}(y),\qquad y\in Y,
\]
and lies in the epigraph of a finite strictly convex function
$g$ on an open convex horizontal domain containing $t(Y)$.
Assume $g(s)=|s|^2$ on $C=t(G)$, the signed graph is represented
on $L^2(G,m)$, and a positive barycentric kernel realizes an
affine-agreeing UCP map. Then $D\ge4$.
In particular this mechanism cannot give an affine-dimension-four
counterexample, with or without differentiability of $g$.
\end{corollary}
\begin{proof}
The continuous surjection $t:G\to C$ sends Haar measure to a
measure of full support: the preimage of a nonempty relatively
open set is nonempty and open, and therefore has positive Haar
measure. This does not require injectivity of the encoding.
Some horizontal coordinate has a nonzero coefficient, and
Lemma~\ref{reg:lem:injective} gives an injective defect for the
actual barycentric map.

If $D\le3$, Theorem~\ref{prox:thm:envelope} gives
$g(t)\ge|t|^2-\dist(t,C)^2$. Containment of the encoded points
therefore implies
\[
 \dist(t(y),C)^2-U(y)+c\mathbf1_{Y_0}(y)\ge0\qquad(y\in Y),
\]
contradicting Theorem~\ref{enc:universal} with $\lambda=1$.
\end{proof}

\begin{remark}\label{prox:rem:scope}
Theorem~\ref{prox:thm:envelope} itself does not assume a Walsh
model, stationarity, or geometric vector norms. It uses one
injective affine defect and full spectral support on the
quadratic graph. Its dimension restriction is $D-1\le2$, the
covering exponent of the auxiliary distance level. Accordingly,
the strictly convex graph obstruction applies in affine dimension
four; the stronger affine-dimension-five exclusion in
Theorem~\ref{enc:universal} applies to the distance-envelope
certificate. These are construction-family obstructions, not
universal dimension lower bounds. They leave open constructions
using a different compression, nongeometric vector norms, other
height identities, or geometry outside the strictly convex graph
setting.
\end{remark}

\section{A rank-one defect cannot produce a counterexample}

\begin{theorem}[Rank-one total quadratic defects are impossible]
\label{rankone:thm}
Let $K\subset\R^d$ be a nonempty compact convex set whose extreme
boundary $X=\operatorname{ex}K$ is compact. Let
$\pi:C(X)\to B(H)$ be a unital $*$-representation and let
$\Psi:C(X)\to B(H)$ be a unital completely positive map agreeing
with $\pi$ on the restrictions of affine functions on $K$.
Write $x_1,\ldots,x_d$ for the coordinate functions and set
\[
 D=\sum_{j=1}^d\bigl(\Psi(x_j^2)-\pi(x_j)^2\bigr).
\]
If $\rank D\le1$, then $\Psi=\pi$. In particular, $D$ cannot have
rank one.
\end{theorem}

\begin{proof}
Take a Stinespring dilation and identify its isometric copy of $H$
with a subspace of the dilation space $H\oplus E$. Write
\[
 N_j=\rho(x_j)=
 \begin{pmatrix}A_j&B_j\\ B_j^*&C_j\end{pmatrix},
 \qquad A_j=\pi(x_j).
\]
The $N_j$ commute, the $A_j$ commute, and
\[
 \Psi(x_j^2)-A_j^2=B_jB_j^*,
 \qquad D=\sum_jB_jB_j^*.
\]
Suppose $\rank D=1$ and let $h$ be a unit vector spanning its range.
There are vectors $k_j\in E$ such that $B_j=hk_j^*$, with at least
one $k_j\ne0$. Put $Q=I_H-hh^*$ and $v_j=QA_jh$.
Applying $Q$ to the upper-right block of $N_iN_j=N_jN_i$ gives
\begin{equation}
                 v_i k_j^*=v_j k_i^*.
\label{rankone:eq:cross}
\end{equation}

If $\operatorname{span}_{\C}\{k_1,\ldots,k_d\}$ has dimension at
least two, choose independent $k_i,k_j$. Equation
\eqref{rankone:eq:cross} implies $v_i=v_j=0$: evaluate it on a
vector orthogonal to one of $k_i,k_j$ but not the other.
Pairing any remaining index with $i$ then gives $v_\ell=0$.
Thus $h$ is a joint eigenvector of the $A_j$.

If that span has dimension one, the upper-left commutation
equations give $B_iB_j^*=B_jB_i^*$. Choosing a nonzero $B_1$,
after relabeling, it follows that
$B_j=\lambda_jB_1$ with $\lambda_j\in\R$ and $\lambda_1=1$.
Equation \eqref{rankone:eq:cross} therefore yields
\[
       (A_j-\lambda_jA_1)h=c_jh,
       \qquad c_j\in\R.
\]
The joint spectral measure of $h$ is supported on the affine line
defined by $x_j-\lambda_jx_1=c_j$ for $j=2,\ldots,d$.
An affine line meets $\operatorname{ex}K$ in at most two points.
Consequently, in both cases $h$ lies in a sum of at most two joint
eigenspaces.

Let $e\ne0$ be a joint eigenvector at $x\in X$. The scalar state
\[
       f\longmapsto
       \frac{\langle\Psi(f)e,e\rangle}{\|e\|^2}
\]
has barycenter $x$. Since $x$ is extreme, its representing
probability measure is $\delta_x$. Hence
\[
 \|B_j^*e\|^2
 =\langle(\Psi(x_j^2)-A_j^2)e,e\rangle=0
 \qquad(1\le j\le d).
\]
Apply this to the at most two spectral components of $h$.
It follows that $B_j^*h=0$ for every $j$, contradicting
$B_j^*h=k_j$ and the existence of a nonzero $k_j$.
This excludes rank one.

If $D=0$, every positive summand $B_jB_j^*$ is zero. Thus all
coordinates belong to the multiplicative domain of $\Psi$.
They generate $C(X)$, so $\Psi=\pi$.
\end{proof}

The theorem concerns the total quadratic defect, rather than Kraus
rank. Its rank is unchanged by an invertible real affine change of
coordinates: the closure of its range is the closed span of the
ranges of the $B_j$. No conclusion for defects of rank two or higher
is asserted here.
\section{A flat-family dimension bound}

\begin{lemma}\label{new:flat}
Let $E$ be a complex Hilbert space, and let $\mathcal U$ be a
finite-dimensional real linear subspace of $B(\C^m,E)$, with $m<\infty$.
If
\[
                         AB^*=BA^*\qquad(A,B\in\mathcal U),
\]
then $\dim_{\R}\mathcal U\le m$.
\end{lemma}
\begin{proof}
For $x\in E$ define the real-linear map
$L_x:\mathcal U\to\C^m$ by $L_x(A)=A^*x$. Choose $x$ maximizing its
real rank $r$. We claim that
\[
                     L_y(\ker L_x)\subset\ran_{\R}L_x
                         \qquad(y\in E).
\]
Indeed, choose real domain and range decompositions in which $L_x$ has
an invertible $r\times r$ upper-left block and zero remaining blocks.
If the lower-right block of $L_y$ were nonzero, an
$(r+1)\times(r+1)$ minor of $L_x+tL_y$ would have a nonzero linear term
in real $t$. Its rank would exceed $r$ for small $t\ne0$, contradicting
$L_x+tL_y=L_{x+ty}$ and the choice of $x$.

Now suppose $A^*x=0$. For every $B\in\mathcal U$,
$AB^*x=BA^*x=0$. Hence $A^*y$ is complex-orthogonal to every $B^*x$,
for every $y\in E$. The claimed inclusion puts $A^*y$ in the real span
of those same vectors. Therefore $A^*y=0$ for every $y$, so $A=0$.
Thus $L_x$ is injective.

The real subspace $\ran_{\R}L_x\subset\C^m$ is isotropic for
$(v,w)\mapsto\operatorname{Im}\langle v,w\rangle$: the cross-product
hypothesis makes $AB^*$ self-adjoint, so
$\langle A^*x,B^*x\rangle$ is real. This subspace is real-orthogonal to
its product by $i$; its real dimension is at most $m$. Injectivity
concludes the proof.
\end{proof}

This resembles the flat-bilinear-form identities from submanifold
geometry, but the proof requires no classification theorem and handles
rank-deficient families.

\section{Block Hardy compressions}

Let $T_F=P_+M_F|_{H^2\otimes\C^m}$ for a bounded measurable
$m\times m$ matrix symbol $F$. Let $S$ be the multiplicity-$m$ unilateral
shift and $E_0:\C^m\to H^2\otimes\C^m$ the constant-function inclusion. We use
$H_G=(I-P_+)M_G|_{H^2\otimes\C^m}$ for the Hankel operator.

\begin{proposition}\label{new:block}
A finite commuting self-adjoint block Toeplitz tuple of multiplicity $m$
admits a finite orthogonal decomposition by constant coefficient spaces
\[
                         \C^m=\bigoplus_\alpha E_\alpha
\]
such that on each block both its joint spectrum and the
multiplication-symbol joint spectrum lie in the same affine subspace of
dimension at most $\dim_\C E_\alpha$. The coefficient dimensions sum
to $m$.
\end{proposition}
\begin{proof}
For completeness, commuting Toeplitz operators have commuting symbols.
The semicommutator identity
$T_{FG}-T_FT_G=H_{F^*}^*H_G$ and
$H_GS^n\to0$ strongly imply
$S^{*n}T_FT_GS^n\to T_{FG}$ strongly. Here the convergence of the
Hankel tails follows first on analytic polynomials from square summability
of Fourier coefficients and then by boundedness. Thus
$[T_F,T_G]=0$ implies $T_{FG-GF}=0$, hence $FG=GF$ almost everywhere,
by injectivity of the Toeplitz-symbol map. Self-adjoint Toeplitz operators
likewise have Hermitian symbols.

Write $\mathcal W$ for the real span of the identity and the tuple.
For $A\in\mathcal W$ put $U_A=S^*AE_0$. Since $S^*AS=A$,
\[
 S^*[A,B]S-[A,B]=U_AU_B^*-U_BU_A^*.
\]
Lemma~\ref{new:flat} bounds the real dimension of
$\{U_A:A\in\mathcal W\}$ by $m$. The kernel of $A\mapsto U_A$
consists exactly of constant Hermitian matrices $I\otimes C$, since
$U_A$ records every positive Fourier coefficient and self-adjointness
determines the negative coefficients.

If this kernel consists only of scalar constants, then
$\dim_\R\mathcal W\le m+1$. The coordinate joint spectrum therefore
lies in an affine subspace of dimension at most $m$. Every affine
relation is also a relation between the symbols, by Toeplitz
injectivity, so their joint spectrum lies in the same subspace.

Otherwise choose a nonscalar constant Hermitian matrix in the kernel.
It commutes with the whole tuple. Its eigenspaces reduce both the
Toeplitz operators and their symbols: the corresponding symbol
commutators have zero Toeplitz operator and hence vanish. Split into
these proper coefficient subspaces and repeat the argument on each
restricted tuple. Each split decreases the relevant positive integer
coefficient dimension, so the procedure terminates in blocks with
only scalar constants in their kernels. The previous paragraph then
applies to each block.
\end{proof}

\begin{corollary}\label{new:hardy}
Let $K\subset\R^d$ be compact convex with $X=\ex K$ compact, and let
$\rho:C(X)\to L^\infty(\mathbb T,M_m)$ be a unital representation.
Suppose compression to $H^2\otimes\C^m$ agrees on the affine coordinates
with a representation $\pi$ of $C(X)$. If $m\le3$, the compression is $\pi$ on all of $C(X)$.
\end{corollary}
\begin{proof}
Apply Proposition~\ref{new:block} to the compressed coordinate tuple.
The constant coefficient blocks reduce the full representations,
because the coordinates generate $C(X)$. On each block let $L$ be the
common affine subspace. Both representations factor through
$C(Y)$, where $Y=X\cap L$ is compact. Moreover
$Y=\ex(\operatorname{conv}Y)$: each point of $Y$ remains extreme in
the smaller convex set, and extreme points of the convex hull of a
compact set belong to that set. Ambient affine functions restrict onto
all affine functions of this convex hull.

For $m\le2$, the planar theorem applies on every block. For $m\le3$,
Theorem~\ref{thm:three} covers the only additional case,
after identifying the affine span with $\R^3$. Each compression is
therefore the source representation, and so is their finite sum.
\end{proof}

\begin{remark}
The target-symbol spectrum must lie in the same affine subspace as the
source spectrum; a statement about the source alone would not justify
the planar reduction. Toeplitz injectivity supplies precisely this
additional fact. No regularity of the symbols is assumed. Arbitrary
infinite-dimensional compression subspaces are not covered.
\end{remark}

\section{Finite POVMs with a persistent quadratic defect}

Put $Q(x)=\sum_{j=1}^d x_j^2$. The next equivalence gives a route through
finite semidefinite feasibility problems that does not prescribe a Walsh
dilation, geometric decay, or an additive height correction.

\begin{proposition}\label{new:finite}
For fixed compact convex $K\subset\R^d$ with compact $X=\ex K$,
$A(K)|_X$ is not hyperrigid if and only if there are $\delta>0$ and,
for arbitrarily small $\epsilon>0$, finite source and target lists
$x_1,\ldots,x_N,y_1,\ldots,y_M\in X$, positive matrices
$E_1,\ldots,E_M\in M_N$, and a unit vector $h\in\C^N$ satisfying
\[
 \sum_kE_k=I_N,\qquad
 \left\|\sum_k(y_k)_jE_k-\diag((x_i)_j)\right\|\le\epsilon
                           \quad(1\le j\le d),
\]
and
\[
 \left\langle\left[\sum_kQ(y_k)E_k-\diag(Q(x_i))\right]h,h\right\rangle
 \ge\delta.
\]
Repeated source points are allowed. Both finite list lengths may vary,
but $K,X$, and $\delta$ are fixed.
\end{proposition}
\begin{proof}
First take a sequence of finite data with $\epsilon_n\to0$. Define
\[
 \pi_n(f)=\diag(f(x_i)),\qquad \Phi_n(f)=\sum_kf(y_k)E_k.
\]
These are a representation and a UCP map. Their classes define maps
into the quotient C*-algebra
\[
             \mathcal A=\prod_nM_{N_n}\big/\bigoplus_nM_{N_n}.
\]
The first remains a representation; the second is UCP. They agree on
every affine coordinate, while the quadratic difference has quotient
norm at least $\delta$. A faithful representation of $\mathcal A$ on
a Hilbert space gives a nontrivial agreeing pair and hence failure of
hyperrigidity. If a separable space is desired, the orbit of a witness
under the maps on a countable dense family and their adjoints gives a
separable reducing subspace that retains the discrepancy.

Conversely let $\pi,\Phi$ be a nontrivial agreeing pair. The coordinate
defect identity implies
\[
             \Phi(Q)-\pi(Q)=\sum_j
             [\Phi(x_j^2)-\pi(x_j)^2]\ge0,
\]
and this operator is nonzero, since simultaneous vanishing makes the
Stinespring compression multiplicative on the coordinate generators.
Choose a unit vector $h$ with positive expectation $\delta_0$.

Let $P$ be the spectral measure of $\pi$. Partition $X$ into finitely
many Borel cells $B_i$ of diameter at most $\eta$, choose $x_i\in B_i$,
and form the finite-dimensional span of the nonzero vectors $P(B_i)h$.
It contains $h$. The step multipliers $\sum_i f(x_i)P(B_i)$ preserve
this span and give an exact finite-dimensional representation. Compress
$\Phi$ to the same space. Its coordinate differences from the step
representation are at most $\eta$, while its quadratic discrepancy at
$h$ is at least $\delta_0-\omega_Q(\eta)$, where $\omega_Q$ is the
uniform modulus of continuity of $Q$ on $X$.

Partition the target $X$ into small Borel cells $C_k$, choose
$y_k\in C_k$, and replace each effect of the compressed POVM on $C_k$
by an atom at $y_k$. The effects are positive and sum to the identity.
Positivity bounds the coordinate changes by the mesh and the quadratic
change by its modulus of continuity. Let both meshes tend to zero and
fix $0<\delta<\delta_0$.
\end{proof}

For fixed point lists and $h$, the displayed constraints are linear
matrix inequalities, and the quadratic expectation is a linear objective
in the effects. A single positive numerical result at nonzero
$\epsilon$ is not sufficient. One must prove a uniform positive defect
as $\epsilon\downarrow0$ on a fixed compact extreme boundary. The
matrix sizes necessarily become unbounded: bounded source sizes would,
by compactness of the source points, unit vectors, and finite-matrix
UCP maps, yield an exact finite-dimensional witness, contrary to the
unique extension property of atomic representations.

\section{Connections with other construction methods}
\label{sec:literature}

The relevant literature spans convex geometry and normal-cone stratification,
operator synthesis, potential theory, matrix convexity and semidefinite
interpolation, Hardy-space compressions, and quantum-channel fixed
points. The following connections have specific consequences for this
problem. Results quoted here are distinguished from the arguments proved
in this note. The comparison does not constitute a complete novelty audit.

\subsection{Convex geometry and operator synthesis}

Menne--Santilli's stratification theorem relates the number of independent
normal directions to rectifiable strata \cite[Theorem 4.12]{MS}.
This suggests separating smooth points, edge points, and vertices.
One must retain exceptional sets: an arbitrary spectral measure can be
concentrated on a set of zero surface measure. The countability argument
in Lemma~\ref{lem:countable} supplies the stronger, exact statement needed
here.

Shulman--Turowska connect operator support to linear operator equations
and give a synthesis theorem for coefficient images of dimension at
most two \cite[Theorem 11.1]{ST}. This parallels the two-dimensional normal
lift in the energy argument above. It suggested searching for a way
to preserve the normal label across a dilation. The final proof takes
a simpler route: each remaining source point has at most one nonsmooth contact target,
and the replacement in Proposition~\ref{prop:partner} preserves the
entire UCP map. No general operator-synthesis theorem or bounded
decomposition over a finite-incidence relation is assumed.

\subsection{Potential theory and noncommutative capacity}

Voiculescu relates the quasicentral modulus to nonlinear condenser
capacity, with a precise correspondence for shifts on Cayley graphs
\cite{Voi}. This suggests treating a dilation defect as a quantity that
must survive at arbitrarily small scales. The analogy is useful, but
positive scalar capacity alone cannot provide a dilation.

There is an elementary warning already in dimension four. Surface
measure on $S^3$ has finite inverse-square energy: near a point, the
integral is bounded by a constant times
$\int_0^1 r^{-2}r^2\,dr$. Nevertheless its coordinate system is
hyperrigid. If an agreeing UCP map is defined on $C(S^3)$, then
\[
 \sum_{j=1}^4\bigl(\Phi(x_j^2)-\pi(x_j)^2\bigr)
       =\Phi(1)-I=0.
\]
Each summand is positive, so every coordinate is in the multiplicative
domain. Thus a capacity-based construction must solve joint positivity
and moment identities, not only find a measure with finite energy.

\subsection{Choquet order, matrix convexity, and finite interpolation}

Davidson--Kennedy identify the unique extension property of a
multiplication representation with maximality in the dilation order
\cite[Theorem 7.6]{DK}. Their dilation order differs from classical
Choquet order; the distinction is exactly why ordinary barycentric
transport is insufficient here. At an extreme point a classical
barycentric probability kernel is necessarily a point mass.

The matrix-range framework of Davidson--Dor-On--Shalit--Solel
\cite{DDSS} and the minimal and maximal matrix convex sets of
Passer--Shalit--Solel \cite{PSS} suggest a different construction route:
positive matrix-valued interpolation on increasingly fine point sets.
Proposition~\ref{new:finite} makes the required limiting statement
precise. General scaled normal-dilation theorems do not themselves
preserve an extreme represented spectrum. Similarly, a matrix convex
gap at a fixed matrix size does not by itself give our diffuse
counterexample.

\subsection{Quantum information and fixed observables}

Quantum-channel theory distinguishes observables fixed by a channel
from those commuting with its operation elements. Magajna proves that
a countable family of commuting normal Kraus operators on a separable
Hilbert space forces the latter conclusion for a normal channel
\cite[Theorem 4.6]{Mag}; Carbone--Jen\v{c}ov\'a study fixed-point algebras
under a faithful normal invariant state \cite[Proposition 4]{CJ}.
Two elementary consequences help filter proposed constructions.

First, if a UCP channel $T$ fixes commuting self-adjoint operators $A_j$
and has a faithful invariant state $\omega$, then
\[
 T(A_j^2)-A_j^2\ge0,\qquad
 \omega\bigl(T(A_j^2)-A_j^2\bigr)=0.
\]
Faithfulness forces equality. The multiplicative-domain argument then
shows that $T$ fixes $C^*(A_1,\ldots,A_d)$. This proof does not require
finite-dimensional Hilbert space. Second, within that normal-channel setting, commuting normal Kraus
operators are excluded by Magajna's theorem. A channel construction
therefore needs operation elements outside this commuting class and
cannot have a faithful invariant state. These conditions are necessary,
not a construction.

\subsection{Hardy spaces and finite multiplicity}

The Brown--Halmos displacement identity and the matrix-symbol
Toeplitz/Hankel framework of Gu--Zheng \cite{BH,GZ} make positivity of
a compressed representation automatic. Exact commutativity of the
compressed coordinates is a separate constraint. The flat-family
argument in this note shows that simply adding two or three internal
components to scalar Hardy space cannot work. Multiplicity at least
four is an unexcluded case of this mechanism, not evidence of an
example. Smooth or finite-band symbols may bring further restrictions
that the present argument does not decide.

\subsection{Dimension and range restrictions}

Tian reduces finite-dimensional extension problems to three-dimensional
operator systems by finite matrix amplification \cite{Tian}. The
generated algebra becomes $M_{3r}(A)$; it is noncommutative even when
$A$ is commutative. Thus its three is a complex operator-system
dimension and is not the affine dimension sought here. The tight
extension results of Clou\^atre--Thompson \cite{CT} restrict the range
to the represented von Neumann algebra and do not exclude arbitrary
$B(H)$-valued extensions. Scherer's spectrahedral theorem
\cite[Theorem 3.7]{spectra} applies to compact spectrahedra defined by a finite
matrix pencil with closed extreme boundary; no extension to their
projections is used.

\section{A nonspectrahedral four-dimensional positive example}
\label{sec:testbody}

A mixed cubic graph initially seemed to offer a simple test beyond the
Walsh construction. Its geometry is explicit, and it escapes the
spectrahedral theorem. A further positivity calculation nevertheless
rules it out. This example is useful precisely because it separates
those two obstructions.

\begin{proposition}\label{prop:testbody}
For $t=(t_1,t_2,t_3)\in[-1,1]^3$, put
\[
 g(t)=|t|^2+\tfrac14t_1t_2t_3,\qquad
 K=\{(t,z):t\in[-1,1]^3,\ g(t)\le z\le4\}.
\]
Then $K$ is a compact convex body of affine dimension four, and
\[
 \ex K=\{(t,g(t)):t\in[-1,1]^3\}
       \cup\{(v,4):v\in\{-1,1\}^3\}
\]
is closed. The body $K$ is not a spectrahedron.
\end{proposition}
\begin{proof}
The Hessian of $g$ has diagonal entries $2$ and off-diagonal entries
$t_k/4$. The sum of absolute values of the off-diagonal entries in
each row is at most $1/2$, so $\nabla^2g\ge(3/2)I$ on the cube.
In particular $g$ is strongly convex. Also $g\le13/4<4$ on the cube,
so $K$ is compact and has nonempty interior.

Strong convexity makes every point of the lower graph extreme. Every
point strictly between the lower graph and the top is in a nontrivial
vertical segment. A top point is extreme precisely when its horizontal
coordinate is a cube vertex. These observations prove the displayed
formula and closedness.

Suppose $K=\{(t,z):L(t,z)\ge0\}$ for a finite Hermitian affine matrix
pencil. At the interior point $e=(0,0,0,1)$, any vector in
$\ker L(e)$ lies in the common kernel of the pencil: use positivity
at $e\pm\epsilon v$ for each direction $v$. Remove this common
kernel, so that $L(e)>0$. The polynomial $p=\det L$ is nonzero and
vanishes on the open lower-graph patch, because positive definiteness
there would make a boundary point interior. Polynomial division in
$z$ now shows that
\[
                 r(t,z)=z-|t|^2-\tfrac14t_1t_2t_3
                 \quad\hbox{divides }p(t,z).
\]
Indeed, the remainder is a polynomial in $t$ vanishing on an open cube.
Along the line $e+s(1,1,1,-16)$ the factor $r$ is
\[
             1-16s-3s^2-\tfrac14s^3.
\]
Its derivative is $-\tfrac34(s+4)^2-4<0$, so the cubic has exactly
one real root, which is simple, and two nonreal roots. In contrast,
\[
 \det L(e+sv)=\det L(e)\,
 \det\bigl(I+sL(e)^{-1/2}L(v)_{\mathrm{lin}}L(e)^{-1/2}\bigr)
\]
has only real zeros, because the middle matrix coefficient is
Hermitian. Divisibility is a contradiction.
\end{proof}

This test also avoids the thin-source obstruction. Restrict the lower
graph to an interior cube, and push forward three-dimensional Lebesgue
measure under its normal lift. The lift is Lipschitz there and its
inverse, metric projection followed by horizontal projection, is
Lipschitz. The pushed-forward measure has finite inverse-square energy,
by comparison with $\int_0^1r^{-2}r^2\,dr$. This only shows that the
earlier zero-capacity test does not apply.

\subsection{Why the cubic perturbation is still rigid}

\begin{lemma}\label{pert:word}
Let $N_1,\ldots,N_d$ be self-adjoint contractions on $H\oplus L$, with
blocks
\[
 N_i=\begin{pmatrix}A_i&B_i\\B_i^*&C_i\end{pmatrix},
 \qquad D=\sum_{i=1}^d B_iB_i^*.
\]
For any word $i_1\cdots i_k$, writing $P$ for the projection onto $H$,
\[
 \left\|PN_{i_1}\cdots N_{i_k}|_H-A_{i_1}\cdots A_{i_k}\right\|
 \leq \binom{k}{2}\|D\|.
\]
The right-hand side is understood as zero for $k=0,1$.
\end{lemma}
\begin{proof}
Put $d_0=\|D\|^{1/2}$. Each $\|B_i\|\leq d_0$, and
$\|[P,N_i]\|=\|B_i\|$. The commutator product rule therefore gives
\[
 \|PN_{i_1}\cdots N_{i_r}(I-P)\|\leq r d_0.
\]
In multiplying a word of length $k-1$ by $N_{i_k}$, its upper-left
block becomes its previous upper-left block times $A_{i_k}$, plus
its previous upper-right block times $B_{i_k}^*$. All $A_i$ are
contractions. Thus the error bound increases by at most
$(k-1)d_0^2$. Induction proves the assertion.
\end{proof}

For a real polynomial $p(t)=\sum_\alpha c_\alpha t^\alpha$ define
\[
             C(p)=\sum_{|\alpha|\geq2}
                      |c_\alpha|\binom{|\alpha|}{2}.
\]
Constant and linear terms contribute nothing to this quantity.

\begin{theorem}\label{pert:graph}
Let $X$ be a compact subset of $[-1,1]^d\times\R$, with coordinate
functions $(t_1,\ldots,t_d,z)$. Suppose
\[
       z\geq g(t):=\sum_{i=1}^d t_i^2+p(t)\quad\hbox{on }X,
       \qquad C(p)<1.
\]
Every unital representation $\pi:C(X)\to B(H)$ supported on
$G=\{(t,z)\in X:z=g(t)\}$ has the unique extension property relative
to the affine coordinate system. No convexity hypothesis on $g$ is
needed for this assertion.
\end{theorem}
\begin{proof}
Let $\Phi:C(X)\to B(H)$ be unital completely positive and agree with
$\pi$ on the coordinate functions. Take a Stinespring dilation
$\Phi(f)=V^*\rho(f)V$, and identify $VH$ with $H$.
The commuting self-adjoint contractions $N_i=\rho(t_i)$ have the
blocks in Lemma~\ref{pert:word}, with $A_i=\pi(t_i)$.
In particular the $A_i$ commute. Put
\[
 E=\Phi(p(t))-p(A_1,\ldots,A_d).
\]
Applying that lemma to a fixed ordering of every monomial gives
$\|E\|\leq C(p)\|D\|$. Since $\pi$ is supported on $G$,
$\Phi(z)=\pi(z)=g(A)$. Positivity of $z-g(t)$ consequently yields
\[
       0\leq\Phi(z-g(t))=-D-E,
       \qquad 0\leq D\leq-E.
\]
Here $E$ is self-adjoint. Hence
$\|D\|\leq\|E\|\leq C(p)\|D\|$, forcing $D=0$ and every $B_i=0$.
The horizontal coordinates therefore intertwine:
$\rho(t_i)V=V\pi(t_i)$, and the same holds for their polynomial $g$.

Set $h=z-g(t)\geq0$. Now $\Phi(h)=0$, so
$\rho(h)^{1/2}V=0$ and hence $\rho(h)V=0$. It follows that
\[
          \rho(z)V=\rho(g(t))V=Vg(A)=V\pi(z).
\]
The real coordinate functions generate $C(X)$. Their intertwining
therefore proves $\Phi=\pi$ on all of $C(X)$.
\end{proof}

\begin{corollary}\label{pert:cubic}
For $|\epsilon|<1/3$, let
\[
 g_\epsilon(t)=t_1^2+t_2^2+t_3^2+\epsilon t_1t_2t_3,
 \qquad
 K_\epsilon=\{(t,z):t\in[-1,1]^3,\ g_\epsilon(t)\leq z\leq4\}.
\]
Then $K_\epsilon$ is a four-dimensional compact convex body with
closed extreme boundary, and its affine coordinate system is
hyperrigid on $C(\ex K_\epsilon)$. In particular this holds for the
previously proposed value $\epsilon=1/4$.
\end{corollary}
\begin{proof}
The Hessian of $g_\epsilon$ has diagonal entries $2$ and off-diagonal
row sums bounded by $2|\epsilon|$. Thus
$\nabla^2g_\epsilon\geq(2-2|\epsilon|)I>0$ on the cube.
Also $g_\epsilon\leq3+|\epsilon|<4$. The body is compact, convex,
and has nonempty interior, and
\[
 \ex K_\epsilon=
 \{(t,g_\epsilon(t)):t\in[-1,1]^3\}
 \ \cup\ \{(\sigma,4):\sigma\in\{-1,1\}^3\}.
\]
Indeed, the lower graph is exposed by the strict supporting tangent
planes of $g_\epsilon$; a point strictly between the lower graph and
the roof lies on a nontrivial vertical segment; and the only extreme
points of the roof are its eight vertices. This extreme boundary is
compact.

For $p(t)=\epsilon t_1t_2t_3$, Theorem~\ref{pert:graph} applies because
$C(p)=3|\epsilon|<1$. Thus every representation supported on the
lower graph has the unique extension property. Each of the eight
remaining evaluations has that property because its point is extreme.
Their arbitrary multiplicities also have the property, and the unique
extension property is preserved by direct sums. Decomposing an
arbitrary representation along the lower graph and the eight roof
points proves the corollary.
\end{proof}

For the cubic alone, the controlling identity is particularly short:
\[
 \Phi(t_1t_2t_3)-A_1A_2A_3
 =A_1B_2B_3^*+B_1B_2^*A_3+B_1C_2B_3^*.
\]
Each summand has norm at most $\|D\|$. This gives the constant $3$
directly. The result makes the cubic body a positive example beyond
the finite-spectrahedron class, by Proposition~\ref{prop:testbody};
it is not a counterexample candidate.



\section{The remaining low-dimensional questions}
\label{sec:questions}

The lower-dimensional obstruction is now complete:
\[
                         4\le d_{\min}\le33.
\]
No construction in affine dimension four is obtained here. The signing
and proximal arguments exclude direct modifications of the stationary
geometric Walsh model; the correlated expectation obstruction and the
new rank-one, Hardy, and polynomial results exclude other simple
mechanisms. These are restrictions on particular constructions, not
evidence that all four-dimensional bodies are hyperrigid.

Proposition~\ref{new:finite} suggests a different route: fix a compact
extreme boundary in dimension four and seek finite POVMs whose affine
operator-norm errors vanish while a positive quadratic discrepancy
persists. Both the boundary and the positive discrepancy must stay
fixed as matrix sizes grow. The cubic body in
Proposition~\ref{prop:testbody} is excluded by
Corollary~\ref{pert:cubic}, despite its nonspectrahedral geometry.
A different body must escape that polynomial estimate.

A channel construction must likewise avoid a faithful invariant state
and the commuting normal Kraus class discussed above. Hardy compression
has first unexcluded multiplicity four. In either approach, exact
commutation, extremality of the represented spectrum, and a nonzero
coordinate defect must all be verified; none follows from the other two.

\subsection*{Preparation and verification}
This manuscript was developed and edited with substantial language-model
assistance. Separate automated review tasks rederived the compression
identities, checked the geometric and measure-theoretic arguments,
and tested proposed constructions against the stated hypotheses. In
version two, a fresh adversarial review rederived the dilation
replacement and replaced the residual disintegration with bounded Borel
endpoint interpolation. The source decomposition, selectors, and
nonseparable reduction were checked separately.
Exact rational arithmetic verifies the numerical certificate in
Section~\ref{sec:construction}. These checks are not independent human
mathematical verification and are not a formal proof certification.
The bibliography identifies the source versions used; the limited
literature comparison does not establish priority for every subsidiary
lemma. The presentation separates complete arguments from the open
questions above.

\appendix
\section{Normal-lift capacity is not a necessary condition}
\label{app:capacity}

The normal lift gives a convenient metric for the supporting-gap
criterion, but it need not give the strongest one. The next example
separates the two tests. Its represented source has the unique extension
property even though its normal lift supports a measure of finite
inverse-square energy.

\begin{theorem}\label{geom:thm}
There are a compact convex body $K\subset\R^4$ with closed extreme
boundary and a compact set $Z\subset\ex K$ such that:
\begin{enumerate}
\item every point of $Z$ is exposed and has a unique unit outward normal;
\item the normal lift of $Z$ supports a probability measure of finite
inverse-square energy;
\item every representation concentrated on $Z$ has UEP relative to $A(K)|_{\ex K}$.
\end{enumerate}
\end{theorem}
\begin{proof}
We construct a thin Cantor set whose normal field records two additional
independent collections of digits. For signs $\omega_n\in\{-1,1\}$,
$n\ge0$, put
\[
 s(\omega)=\left(\sum_{n\ge0}4^{-n}\omega_n,
       10^{-3}\sum_{n\ge0}64^{-n}\omega_n,
       10^{-3}\sum_{n\ge0}128^{-n}\omega_n\right),
 \qquad C=s(\{-1,1\}^{\N_0}).
\]
The first coordinate is an injective Cantor expansion, so $s$ is a
homeomorphism onto $C$. Its first three vector coefficients are
linearly independent by the Vandermonde determinant; hence
$\aff C=\R^3$. Also $C\subset(-2,2)^3$ and $|s|<2$.

If two sequences first differ at index $n$, their first-coordinate
difference is at least $(4/3)4^{-n}$, whereas their second- and
third-coordinate differences are at most
$2\cdot10^{-3}64^{-n}/(1-1/64)$ and
$2\cdot10^{-3}128^{-n}/(1-1/128)$, respectively. In particular,
\begin{equation}\label{app:eq:flat}
 |s_j-t_j|\le\frac{8}{7000}|s_1-t_1|^2
       \quad(j=2,3),\qquad |s-t|^2\le2|s_1-t_1|^2.
\end{equation}
Define continuous functions $a,b:C\to[1,2]$ by
\[
 a(s(\omega))=1+\sum_{k\ge0}2^{-(k+1)}\frac{1+\omega_{3k+1}}2,
 \qquad
 b(s(\omega))=1+\sum_{k\ge0}2^{-(k+1)}\frac{1+\omega_{3k+2}}2.
\]
Set $p(s)=2s+(0,a(s),b(s))$. The preceding estimates imply
\begin{equation}\label{app:eq:interpolation}
 \bigl|\langle p(s)-2s,t-s\rangle\bigr|
       \le\frac{32}{7000}|t-s|^2<\frac12|t-s|^2
       \qquad(s,t\in C).
\end{equation}
Thus the data $g(s)=|s|^2$, $\nabla g(s)=p(s)$ admit the strongly
convex extension
\[
 g(x)=\max_{s\in C}
       \left\{|s|^2+\langle p(s),x-s\rangle+\tfrac12|x-s|^2\right\}.
\]
Indeed, subtracting $|x|^2/2$ leaves a supremum of affine functions.
For $x\in C$, \eqref{app:eq:interpolation} shows that the unique
maximizer is $s=x$, and $g(x)=|x|^2$.

For clarity, this uniqueness also proves differentiability at every
$x\in C$. If $s_h$ maximizes at $x+h$, compactness and continuity
force $s_h\to x$. Let $Q_s$ denote the quadratic inside the maximum.
Then
\[
 Q_x(x+h)-Q_x(x)\le g(x+h)-g(x)
      \le Q_{s_h}(x+h)-Q_{s_h}(x).
\]
Both bounds are $\langle p(x),h\rangle+o(|h|)$.

Choose $M>\max_{[-2,2]^3}g$ and let
\[
 K=\{(x,z):x\in[-2,2]^3,\ g(x)\le z\le M\},
 \qquad Z=\{P_s=(s,|s|^2):s\in C\}.
\]
Strict convexity shows that the full extreme boundary consists of the
lower graph and the eight vertices of the top face, exactly as in
Lemma~\ref{lem:extreme}. It is compact, and $K$ has nonempty interior.
At $P_s$, the supporting plane is strict away from $P_s$, and the unique
outward unit normal is
\[
                         n(s)=\frac{(p(s),-1)}{\sqrt{1+|p(s)|^2}}.
\]

We next show that the lift $P_s+n(s)$ carries finite energy. The paired
lift $(P_s,n(s))$ and the sum lift $P_s+n(s)$ are bi-Lipschitz equivalent:
the supporting-normal inequalities give
\[
 |P_s-P_t|^2+|n(s)-n(t)|^2
 \le |P_s+n(s)-P_t-n(t)|^2
 \le2\bigl(|P_s-P_t|^2+|n(s)-n(t)|^2\bigr).
\]
By \eqref{app:eq:flat}, the paired lift is in turn bi-Lipschitz equivalent
to the coordinates $(s_1,a(s),b(s))$. In one direction all its components
are Lipschitz functions of these three coordinates. In the other,
recover $s$ from the position and then recover $p(s)$ from the normal
by division by its last coordinate, which is bounded away from zero.

Under product measure on the signs, $a$ and $b$ are independent uniform
variables on $[1,2]$. Conditional on the digits defining them, the
remaining contribution to $s_1$ is a translate of
\[
                         \sum_{k\ge0}64^{-k}\omega_{3k}.
\]
Its probability in an interval of length $2r$ is bounded by a fixed
multiple of $r^{1/6}$. To verify this, choose $N$ with
$64^{-N}\le r<64^{-(N-1)}$. An interval of that length meets at most
two level-$N$ cylinders, each of mass $2^{-N}$; the separation follows
by comparing the first differing digit with the geometric tail.
Conditioning and integrating over $a,b$ therefore give a ball bound
of the form $Cr^{2+1/6}$ for $(s_1,a,b)$, uniformly in the center.
The same bound holds on the normal lift by bi-Lipschitz equivalence.
For its pushed-forward probability measure $\nu$, it follows that
\[
 \iint |u-v|^{-2}\,d\nu(u)d\nu(v)
 \le1+2C\int_0^1r^{-3}r^{13/6}\,dr<\infty.
\]

Finally, consider the actual supporting affine functions
\[
 b_{P_s}(x,z)=z-|s|^2-\langle p(s),x-s\rangle.
\]
They expose $P_s$, since
$b_{P_s}(x,z)\ge z-g(x)+|x-s|^2/2$ on $K$. On the represented source,
\eqref{app:eq:flat}--\eqref{app:eq:interpolation} give
\[
                  0\le b_{P_s}(P_t)\le3|s_1-t_1|^2.
\]
The first coordinate lies in the compact interval $[-4/3,4/3]$,
which supports no nonzero measure of finite inverse-square energy.
Theorem~\ref{bar:thm} proves the unique extension property on $Z$.
\end{proof}

\section{Varying distributions and uniformly positive defects}
\label{app:heterogeneous}

The local defect estimate in Lemma~\ref{het:lem:local} imposes a
restriction independent of the choice of convex graph. Here is its
consequence for varying three-point distributions.

At coordinate $n$, let the symmetric distribution on $\{0,\pm R_n\}$
have second moment $b_n$, where $1<b_n<R_n^2$. Put $a_n=b_n^{-1}$.
Use the branch spaces of the Walsh construction, with these varying
tail distributions. Replace the formula for $V$ by
\[
 \begin{split}
 V\chi_F={}&\left(\prod_{n\in F}\sqrt{a_n}\right)\chi_F^{(0)}\\
 &+\sum_{k\in F}\sqrt{1-a_k}
     \left(\prod_{n\in F_{>k}}\sqrt{a_n}\right)
       \chi_{F\setminus\{k\}}^{(k)},
 \end{split}
\]
where the normalized branch monomial is
$\chi_F^{(k)}=(\prod_{n\in F_{>k}}b_n^{-1/2})\prod_{n\in F}x_n$.
The proof of Lemma~\ref{lem:walsh} remains valid because
\[
 \prod_{n\in F}a_n+
 \sum_{k\in F}(1-a_k)\prod_{n\in F_{>k}}a_n=1.
\]
In particular all square-free products compress exactly. For the
coordinate-square defect, parity and the same branch calculation give
\begin{equation}\label{het:eq:Dn}
 \bigl(\Phi(x_n^2)-I\bigr)\chi_F
 =\left(\prod_{i\in F_{>n}}a_i\right)
 \begin{cases}
 b_n-1,&n\notin F,\\
 R_n^2/b_n-1,&n\in F
 \end{cases}
 \chi_F.
\end{equation}

\begin{theorem}\label{het:thm:obstruction}
Assume $\sum_n(b_n-1)<\infty$. Let $w_n\in\R^D$ satisfy
$\sum_n R_n|w_n|<\infty$, with at least one $w_n\ne0$, and put
\[
                       t=\sum_n w_nx_n,
                       \qquad s=\sum_nw_n\varepsilon_n.
\]
There is no compact convex body with compact extreme boundary carrying
an affine-agreeing pair $(\pi,\Psi)$ for which the first $D$ ambient
coordinate functions $\tau_j$ satisfy $\pi(\tau_j)=M_{s_j}$ and
\[
                         \Psi(\tau_j^2)\ge\Phi(t_j^2)
                         \qquad(1\le j\le D).
\]
In particular, no positive barycentric-kernel construction realizes
these data, regardless of its additional coordinates.
\end{theorem}
\begin{proof}
The product $a_*:=\prod_n a_n$ is positive. Equation~\eqref{het:eq:Dn}
gives
\[
 \Phi(x_n^2)-I\ge
 a_*\min\{b_n-1,R_n^2/b_n-1\}I>0.
\]
Uniform convergence and exact compression of mixed products imply
\[
 \Phi(t_j^2)-M_{s_j}^2
       =\sum_n(w_n)_j^2\bigl(\Phi(x_n^2)-I\bigr).
\]
For some $j$, a nonzero summand bounds this defect below by a positive
multiple of $I$. The asserted inequality for $\Psi$ contradicts
Lemma~\ref{het:lem:local}. For a barycentric kernel it follows from
scalar Jensen's inequality for a coordinate square and positivity of
compression.
\end{proof}

The same obstruction applies to finite nonlinear coordinates.

\begin{proposition}\label{het:prop:polynomial}
Under the preceding summability hypothesis, every nonconstant real
multiaffine polynomial $p$ in finitely many coordinates satisfies
\[
                    \Phi(p^2)-M_{p(\varepsilon)}^2\ge c_p I
\]
for some $c_p>0$. It therefore cannot be used as an affine coordinate
of a compact-extreme-boundary barycentric construction.
\end{proposition}
\begin{proof}
Exact square-free compression identifies the defect with $T^*T$, where
$T=M_pV-VM_{p(\varepsilon)}$. Restrict this error to the initial branch
$Y_0$. If $p$ involves the first $N$ variables, its tail factor is
bounded below by the positive product $\prod_{n>N}b_n^{-1/2}$.
It suffices to prove injectivity of the finite-prefix error.

Identify functions on the binary cube with multiaffine polynomials.
Let $\mathcal D$ scale a monomial $x_F$ by $\prod_{i\in F}b_i^{-1}$.
If the error vanishes on a vector represented by $h$, then, with
$q=\mathcal Dh$,
\[
                       p(x)q(x)=\mathcal D(p(\varepsilon)h(\varepsilon))(x)
\]
on the full three-point product grid. Both sides have degree at most
two in each variable, so this is a polynomial identity. The right side
is multiaffine. Whenever $p=p_0+x_ip_1$ with $p_1\ne0$, write
$q=q_0+x_iq_1$; comparison of the $x_i^2$ coefficient gives
$p_1q_1=0$, hence $q_1=0$. Thus $q$ and $h$ are independent of every
active variable of $p$. The identity reduces to
\[
        \bigl[p(x)-p(x_1/b_1,\ldots,x_N/b_N)\bigr]q(x)=0.
\]
The first factor is nonzero, because $p$ is nonconstant and every
$b_i>1$. Hence $q=0$. The finite-prefix error has positive minimum
singular value; tensoring with the bounded-below tail proves the claim.
\end{proof}

\section{Correlated magnitudes and an expectation obstruction}
\label{app:correlated}

Independence of the magnitudes is not essential to the Walsh isometry.
This observation gives a distinct family of exact compression maps.
A simple expectation test, however, excludes an entire range of its
possible geometric realizations.

Fix $b>1$ and $0<\delta<b$. Let $\theta$ be uniform on the circle and
let $\xi_n$ be independent uniform signs, independent of $\theta$. Put
\begin{equation}\label{corr:eq:model}
                   x_n=\xi_n\sqrt{b+\delta\cos(4^n\theta)}.
\end{equation}
These coordinates share their angle and are not independent.
Nevertheless, if $S_n=x_n^2$, then for every finite index set and
integers $0\le k_n\le3$,
\begin{equation}\label{corr:eq:moments}
                   \mathbb E\prod_n S_n^{k_n}
                      =\prod_n\mathbb E S_n^{k_n}.
\end{equation}
Indeed, a nonconstant term in the Fourier expansion has frequency
$\sum_n j_n4^n$ with $|j_n|\le3$. Its largest nonzero term cannot be
cancelled, since $4^m>3\sum_{n<m}4^n$. Thus only products of constant
Fourier coefficients survive. By contrast,
\[
 \mathbb E[(S_n-b)^4(S_{n+1}-b)]=\delta^5/16\ne0,
\]
as follows from the $\cos4t$ coefficient in $\cos^4t$.

On the initial branch use \eqref{corr:eq:model} at every coordinate;
on branch $Y_k$ use the binary prefix, the zero-valued star, and the
same correlated law on its tail. The angle and all prefix and tail
signs are independent. For compactness, use disjoint tagged copies of
the binary symbols and of the compact tail intervals, even where their
numerical values agree. Each fixed branch is a compact image of its
angle and sign space; branches with stars tending to infinity
accumulate only on the binary source. Thus the union is compact,
$Y_0$ is clopen, and every numerical coordinate is continuous.

\begin{proposition}\label{corr:prop:compression}
With these branch laws and the same formula for $V$ as in
\eqref{eq:V}, using $a=b^{-1}$, the map $V$ is an isometry and every
square-free product compresses exactly. Moreover $V1=\ind_{Y_0}$ and
\begin{equation}\label{corr:eq:defect}
 \bigl(\Phi(x_n^2)-I\bigr)\chi_F
  =b^{-|F_{>n}|}
  \begin{cases}
   b-1,&n\notin F,\\
   \delta^2/(2b^2),&n\in F
  \end{cases}\chi_F.
\end{equation}
In particular each coordinate-square defect is positive and injective.
\end{proposition}
\begin{proof}
Independent signs make different monomial labels orthogonal.
Equation~\eqref{corr:eq:moments} gives the same normalized monomial
norms as the independent law. A surviving square-free matrix coefficient
uses only exponents zero and two, so the complete proof of
Lemma~\ref{lem:walsh} applies.
For the square defect, $\mathbb E x_n^4=b^2+\delta^2/2$.
If $n\in F$, the total pre-star branch weight is
$b^{-|F_{>n}|-1}$ and the weighted second moment is
$(b^2+\delta^2/2)/b$; the star contributes zero. Subtracting $I$
gives the second coefficient in \eqref{corr:eq:defect}. If $n\notin F$,
the calculation uses the ordinary second moment $b$, giving the first.
\end{proof}

\begin{theorem}[An obstruction independent of the height correction]
\label{avg:lem:vacuum}
For the correlated model, let $w_n\in\R^D$ be a nonzero absolutely
summable sequence and put
\[
 t=\sum_nw_nx_n,\qquad s=\sum_nw_n\varepsilon_n,
 \qquad C=s(G).
\]
Suppose a continuous height $z$ satisfies
\[
 \Phi(z)=M_{|s|^2},\qquad
 z(y)\ge|t(y)|^2-\dist(t(y),C)^2\quad(y\in Y_0).
\]
Then
\begin{equation}\label{avg:eq:mean}
                   \frac1{2\pi}\int_0^{2\pi}
                     \sqrt{b+\delta\cos\theta}\,d\theta\le1.
\end{equation}
No additive form for $z$, and no geometric decay of $|w_n|$, is assumed.
\end{theorem}
\begin{proof}
Let $L=\sum_n|w_n|^2>0$ and denote the mean in
\eqref{avg:eq:mean} by $m$. Exact height compression and
$V1=\ind_{Y_0}$ give $\mathbb E_{Y_0}z=L$, whereas independent signs
give $\mathbb E_{Y_0}|t|^2=bL$.
For each branch point, compare $t$ with the source point having the
same signs. Averaging first over signs and then over the angle gives
\[
 \mathbb E_{Y_0}\dist(t,C)^2
 \le\mathbb E\left|\sum_nw_n\xi_n
           \bigl(\sqrt{b+\delta\cos(4^n\theta)}-1\bigr)\right|^2
 =(b+1-2m)L.
\]
Containment gives the reverse lower bound $(b-1)L$. Dividing by $L$
proves $m\le1$. Absolute uniform convergence justifies all expansions.
\end{proof}

\begin{corollary}\label{avg:cor:first}
The distribution $b=6/5$, $\delta=1$ admits no such exact height in
any finite horizontal dimension. Suppose a body $K$ satisfies the
geometric hypotheses of Theorem~\ref{prox:thm:envelope}, contains
every encoded point $(t(y),z(y))$, and a positive barycentric kernel
realizes the affine-agreeing map on the signed-source representation.
Then no finite strictly convex graph can repair this model when
$D\le3$.
\end{corollary}
\begin{proof}
The two chords of the square-root function on $[b-\delta,b]$ and
$[b,b+\delta]$ give
\[
 m\ge\sqrt b+
  \frac{\sqrt{b+\delta}+\sqrt{b-\delta}-2\sqrt b}{\pi}
 >\frac{\sqrt b+\sqrt{b+\delta}+\sqrt{b-\delta}}3.
\]
The last inequality uses strict concavity and $\pi>3$. At the stated
parameters, the three square roots exceed $109/100$, $148/100$, and
$44/100$, respectively. Thus $m>301/300>1$.
For the graph assertion, exact mixed products and
\eqref{corr:eq:defect} give an injective defect in any nonzero
horizontal coordinate. A positive barycentric kernel preserves that
lower bound, and the signed source has full spectral support.
Theorem~\ref{prox:thm:envelope} forces the excluded quadratic envelope.
\end{proof}

More generally, the mean-radius condition forces
$1<b<\pi^2/8$ and $\delta^2>4(b-1)$. For the first inequality,
the mean is strictly decreasing in $\delta$ and its limit at
$\delta=b$ is $2\sqrt{2b}/\pi$. For the second, its value is strictly
larger than the endpoint chord average
$(\sqrt{b+\delta}+\sqrt{b-\delta})/2$.
These restrictions leave admissible pairs near $(1,1)$; they do not
exclude the whole correlated family or establish containment for any
remaining pair.

\section{Restrictions on direct operator constructions}
\label{app:direct}

We record two short diagnostics that are independent of the Walsh
construction. They explain why two natural direct constructions cannot
supply the desired diffuse example. Neither is a universal dimension
bound.

\begin{proposition}\label{pert:prop:density}
Suppose an affine-agreeing POVM $E$ satisfies $E(B)\le\nu(B)I$ for
every Borel $B\subset X$, where $\nu$ is a finite positive scalar measure.
Then the spectral support of $\pi$ has at most $\lfloor\nu(X)\rfloor$
points and $E=P$. In particular, an operator-valued density with
integrable operator norm cannot yield a diffuse counterexample.
\end{proposition}
\begin{proof}
For a spectral-support point $x$, choose unit vectors
$h_n\in P(B(x,1/n))H$. Their scalar measures for $E$ are probabilities
dominated by $\nu$, with barycenters tending to $x$. Any weak limit is
$\delta_x$ by extremality, and remains dominated by $\nu$. Therefore
$\nu(\{x\})\ge1$. There are at most $\lfloor\nu(X)\rfloor$ such points,
and Lemma~\ref{lem:atomic} gives equality of the maps.
For a positive density $F$ relative to a scalar measure $\mu$, take
$d\nu=\|F\|\,d\mu$.
\end{proof}

This is an order-domination hypothesis, not merely absolute continuity
with respect to a scalar control measure. It does not automatically
apply to a correction of a spectral measure.

\begin{proposition}\label{nonwalsh:prop:toeplitz}
Let $\rho:C(X)\to L^\infty(\mathbb T)$ be a unital
$*$-representation, and let $P_+$ be the orthogonal projection from
$L^2(\mathbb T)$ onto the scalar Hardy space $H^2$. The compression
\[
                   \Phi(f)=P_+M_{\rho(f)}|_{H^2}
\]
cannot be a nontrivial UCP map agreeing with a representation on
$A(K)|_X$.
\end{proposition}
\begin{proof}
Let $S$ be the unilateral shift on $H^2$ and let $e_0$ be its first
basis vector. For self-adjoint Toeplitz operators $A,B$, put
$u_A=S^*Ae_0$ and $u_B=S^*Be_0$. The identities $S^*AS=A$ and
$SS^*=I-e_0e_0^*$ give
\[
                  S^*[A,B]S-[A,B]=u_Au_B^*-u_Bu_A^*.
\]
If $A,B$ commute and $u_A\ne0$, the two rank-one products agree.
It follows that $u_B=\lambda u_A$ for some real $\lambda$.
For a real symbol $f$, $u_{T_f}$ consists of its strictly positive
Fourier coefficients. The negative coefficients are their conjugates.
Consequently, commuting real scalar Toeplitz symbols are all of the
form $\lambda_jf+c_j$ for one real symbol $f$, unless all are constant.

In an affine-agreeing compression the compressed coordinates commute,
because they equal the coordinates of the representation. Their joint
spectral support therefore lies on an affine line, or is a singleton.
A line meets $\ex K$ in at most two points: every interior point of
its intersection segment with $K$ is nonextreme. The represented
support is thus finite, and Lemma~\ref{lem:atomic} applies.
\end{proof}

In contrast, small norm and low Kraus rank do not themselves impose
rigidity. Any counterexample $\Psi(f)=V^*\rho(f)V$ can be amplified
on $H\oplus L\oplus L\oplus\cdots$ to
$\widehat\pi=\pi\oplus\rho\oplus\rho\oplus\cdots$. The isometry
\[
             W(h,\ell_1,\ell_2,\ldots)=(0,Vh,\ell_1,\ell_2,\ldots)
\]
satisfies
$W^*\widehat\pi(f)W=\Psi(f)\oplus\rho(f)\oplus\rho(f)\oplus\cdots$.
Mixing this compression with the identity gives a nontrivial two-Kraus
channel fixing the affine coordinates and arbitrarily close to the
identity. This is a reformulation of a counterexample, rather than an
independent construction of one.

\clearpage
\begin{thebibliography}{DDSS}
\bibitem[Arv]{Arv}
W. Arveson, \emph{The noncommutative Choquet boundary II: Hyperrigidity},
Israel J. Math. \textbf{184} (2011), 349--385.
\url{https://arxiv.org/abs/0810.2751}.

\bibitem[BH]{BH}
A. Brown and P. R. Halmos, \emph{Algebraic properties of Toeplitz operators},
J. Reine Angew. Math. \textbf{213} (1964), 89--102.
\url{https://doi.org/10.1515/crll.1964.213.89}.

\bibitem[BS]{BS}
B. Bilich and M. Scherer, \emph{Commutative counterexamples to Arveson's
hyperrigidity conjecture}, arXiv:2610.26556v1, 2026.
\url{https://arxiv.org/abs/2610.26556v1}.

\bibitem[CJ]{CJ}
R. Carbone and A. Jen\v{c}ov\'a, \emph{On period, cycles and fixed points
of a quantum channel}, arXiv:1905.00857v1, 2019.
\url{https://arxiv.org/abs/1905.00857v1}.

\bibitem[CT]{CT}
R. Clou\^atre and I. Thompson, \emph{Rigidity of operator systems: tight
extensions and noncommutative measurable structures},
arXiv:2406.16806v2, 2025.
\url{https://arxiv.org/abs/2406.16806v2}.

\bibitem[DDSS]{DDSS}
K. R. Davidson, A. Dor-On, O. M. Shalit, and B. Solel,
\emph{Dilations, inclusions of matrix convex sets, and completely positive
maps}, arXiv:1601.07993v3, 2018.
\url{https://arxiv.org/abs/1601.07993v3}.

\bibitem[DK]{DK}
K. R. Davidson and M. Kennedy, \emph{Choquet order and hyperrigidity for
function systems}, arXiv:1608.02334v4, 2021.
\url{https://arxiv.org/abs/1608.02334v4}.

\bibitem[GZ]{GZ}
C. Gu and D. Zheng, \emph{Products of block Toeplitz operators},
Pacific J. Math. \textbf{185} (1998), 115--148.
\url{https://arxiv.org/abs/math/9605225}.

\bibitem[Mag]{Mag}
B. Magajna, \emph{Fixed points of normal completely positive maps on
$B(H)$}, J. Math. Anal. Appl. \textbf{389} (2012), 1291--1302.
\url{https://arxiv.org/abs/1105.1914}.

\bibitem[MS]{MS}
U. Menne and M. Santilli, \emph{A geometric second-order-rectifiable
stratification for closed subsets of Euclidean space},
Ann. Sc. Norm. Super. Pisa Cl. Sci. (5) \textbf{19} (2019), 1185--1198.
\url{https://arxiv.org/abs/1703.09561v2}.

\bibitem[PSS]{PSS}
B. Passer, O. M. Shalit, and B. Solel, \emph{Minimal and maximal matrix
convex sets}, arXiv:1706.05654v2, 2017.
\url{https://arxiv.org/abs/1706.05654v2}.

\bibitem[Sch1]{planar}
M. Scherer, \emph{The hyperrigidity conjecture for compact convex sets
in $\R^2$}, arXiv:2411.11709v1, 2024.
\url{https://arxiv.org/abs/2411.11709v1}.

\bibitem[Sch2]{spectra}
M. Scherer, \emph{The hyperrigidity conjecture for spectrahedra},
arXiv:2601.16075v1, 2026.
\url{https://arxiv.org/abs/2601.16075v1}.

\bibitem[ST]{ST}
V. Shulman and L. Turowska, \emph{Operator synthesis II. Individual
synthesis and linear operator equations}, arXiv:math/0401091v2, 2004.
\url{https://arxiv.org/abs/math/0401091v2}.

\bibitem[Tia]{Tian}
J. Tian, \emph{UCP extension in dimension three},
arXiv:2609.38539v1, 2026.
\url{https://arxiv.org/abs/2609.38539v1}.

\bibitem[Voi]{Voi}
D.-V. Voiculescu, \emph{Capacity and the quasicentral modulus},
arXiv:2107.11924v4, 2022.
\url{https://arxiv.org/abs/2107.11924v4}.
\end{thebibliography}
\end{document}
